\documentclass[10pt]{article}

\usepackage[top=0.85in,bottom=0.85in,left=1.5in,right=1.5in]{geometry}
\usepackage{amsmath,amssymb,amsthm,mathtools}
\usepackage{enumitem}
\setlist[enumerate]{topsep=10pt plus 2pt,partopsep=0pt}
\usepackage{booktabs,longtable}
\usepackage[spaces]{url}
\def\UrlBreaks{\do\.\do\,\do\ \do\*}
\usepackage{xcolor,relsize,xspace}
\usepackage[round]{natbib}
\usepackage[colorlinks,citecolor=blue,linkcolor=black,urlcolor=black]{hyperref}

% ---------------------------------------------------------------------------
% Pointers to the Lean formalization. \leanrepo is the URL of the extracted
% Lean project at the release tag v1.0. \leanref{Part157Defs,Part156Out} cites modules of bsp/lean under
% BstLean, printed as (BstLean:Part157Defs,Part156Out); \leanmod{BSP} cites a
% module outside BstLean. check_coverage.py checks that the groups of the Lean
% files appendix cover every module of the project.
\newcommand{\leanrepo}{https://github.com/andrewcritch/beating-sampling-theorem/tree/v1.0/lean}
\newcommand{\leanmod}[1]{{\relscale{0.78}\color{gray}\hypersetup{urlcolor=gray}%
  \textup{(\href{\leanrepo}{\nolinkurl{#1}})}}}
\newcommand{\leanref}[1]{\leanmod{BstLean:#1}}
\newcommand{\bstlean}{\texttt{BstLean}}
% \leanname{Part289.progOut_run2} names a declaration in namespace BstLean.
\newcommand{\leanname}[1]{\nolinkurl{#1}}
% generated by gen_appendix.py; \leanstepA is S0, \leanstepB is S1, ...
\newcommand{\leanstepA}{Parts 29, 33, 35--36, 41, 48, 53, 56, 58, 73, 80, 105--108, 130--133, 135, 139--140, 142, 144--146, 156--157; CostModel}
\newcommand{\leanstepB}{Parts 141, 150, 157, 159, 169, 201, 272, 278, 289}
\newcommand{\leanstepC}{Parts 134, 137--138, 140, 143, 157, 159, 204--205, 210, 214, 218, 227, 232--233, 236}
\newcommand{\leanstepD}{Parts 105, 107, 109, 112, 134, 136, 148--149, 151--154, 157--158, 160, 188, 194--195, 198--200, 206, 208--209, 211, 215, 217, 219, 221--222, 229, 235, 238, 244--246, 273}
\newcommand{\leanstepE}{Parts 161--171, 202, 207, 212--213, 216, 220, 223, 228, 239, 241--242, 248, 289}
\newcommand{\leanstepF}{Parts 156--157, 192--193, 196--197, 203, 243, 261, 293}
\newcommand{\leanstepG}{Parts 172--188, 190, 278}
\newcommand{\leanstepH}{Parts 243, 249--260, 262--264, 266--271, 278, 285--287}
\newcommand{\leanstepI}{Parts 275--277, 280--281}
\newcommand{\leanstepJ}{Parts 290--293}
\newcommand{\leanstepK}{Parts 150, 189, 279, 282, 288--289}
\newcommand{\leanstepL}{Parts 29, 33, 35--46, 48--53, 56--64, 66, 69, 71--74, 78, 80--82, 84, 86, 88--90, 93, 105--109, 112, 130--135, 138--145, 149--154, 156--180, 182, 185--188, 192--194, 197, 199--201, 204--205, 208--213, 216--217, 220--223, 235--236, 238--239, 241--245, 247, 249, 251--257, 260--262, 264--266, 269, 271, 273, 275--276, 278--279, 281--282, 287--290, 292--293}

\newcommand{\leanparts}[1]{{\relscale{0.78}\color{gray}\hypersetup{urlcolor=gray}%
  \textup{(Lean: \href{\leanrepo}{#1})}}}
% Numbered Lean notes: \leanmark{1} in the text, and after the paragraph
% \begin{leanitems}\leanitem[1]{Part157Paper,Part169Pass}\end{leanitems}.
\newcommand{\leanmark}[1]{{\relscale{0.78}\color{gray}[#1]}}
\newcommand{\leanlist}[1]{\href{\leanrepo}{\nolinkurl{BstLean:#1}}}
% \leanitem[N]{Part1,Part2} lists BstLean modules; \leanitem*[N]{text} links other text to the repository.
\NewDocumentCommand{\leanitem}{s O{1} m}{\item[{[#2]}] \IfBooleanTF{#1}{\href{\leanrepo}{#3}}{\leanlist{#3}}}
\newenvironment{leanitems}{\par\smallskip\relscale{0.78}\color{gray}\hypersetup{urlcolor=gray}\raggedright%
  \begin{list}{}{\setlength{\leftmargin}{2em}\setlength{\labelwidth}{1.6em}%
  \setlength{\labelsep}{0.4em}\setlength{\itemsep}{0pt}\setlength{\topsep}{0pt}%
  \setlength{\parsep}{0pt}}}{\end{list}\par\medskip}
% ---------------------------------------------------------------------------

\emergencystretch=2em
\newtheorem*{bstheorem}{Theorem (Beating Sampling Theorem)}
% \bst links to the Beating Sampling Theorem.
\newcommand{\bst}{\hyperref[thm:main]{Beating Sampling Theorem}}
\newtheorem{lemma}{Lemma}
\newtheorem{conjecture}{Conjecture}
\newcommand{\RR}{\mathbb{R}}
% \Xx is \mathcal{X} and \XX is \mathbb{X} (\AAA, \SSS where \AA, \SS exist).
\newcommand{\Dd}{\mathcal{D}}\newcommand{\Ee}{\mathcal{E}}
\newcommand{\Gg}{\mathcal{G}}\newcommand{\Hh}{\mathcal{H}}
\newcommand{\SSS}{\mathbb{S}}\newcommand{\Kk}{\mathcal{K}}\newcommand{\Nn}{\mathcal{N}}\newcommand{\Bb}{\mathcal{B}}
% Tags for areas to improve: \futurework and \Futurework print the phrase in orange
% (xspace adds a following space where needed).
\newcommand{\futurework}{\textcolor{orange}{future work}\xspace}
\newcommand{\Futurework}{\textcolor{orange}{Future work}\xspace}
\newcommand{\Est}{\mathrm{Est}}
\newcommand{\EE}{\mathbb{E}}
\newcommand{\PP}{\mathbb{P}}
\newcommand{\relu}{\operatorname{ReLU}}
\newcommand{\chat}{\hat c}

\title{Beating Sampling:\\ Mechanistic Estimation of the Failure Rate\\ of a Neural Network at Frontier Scale}
\author{Andrew Critch\\ MAISI}
\date{}

\begin{document}
\maketitle

\begin{abstract}
When can the failure rate of a neural network be estimated more efficiently by analyzing its internals than by sampling its outputs? We investigate this question for a class of random ReLU MLP ``agents'', and prove that the answer is ``usually'', with a single explicit estimation procedure which upper-bounds the failure rate $p$ within a factor of 10, or else by at most 1\%, for at least 99\% of networks and random seeds. The result holds for every input dimension $n\ge8{,}000$. For agents in our class with $8{,}000\le n\le16{,}384$ (comparable to the embedding width of current frontier transformers), the agent MLP has between $4.1\cdot10^{15}$ and $7.2\cdot10^{16}$ parameters, and the procedure costs about as much as a single exact sample of the agent's behavior. By contrast, consider a naive estimator that sees only the outcomes of exact samples and must succeed with probability $0.99$ for every value of the failure rate $p$. When $p \leq 0.1\%$, the naive estimator needs more than $448$ samples on average. For such low failure rates, our procedure is thus more than $447$ times more efficient than naive sampling. The Beating Sampling Theorem and its proof are checked in Lean, and the translation of the Beating Sampling Theorem's assumptions and conclusions into Lean has been checked by GPT-6-Astra (high effort) and Claude Opus~5.5 (medium effort).
\end{abstract}

\section{Introduction}\label{sec:motivation}

This paper is the result of an experiment in human--LLM collaboration on mathematical interpretability research and exposition.  While I, Andrew Critch, wrote the title and this introduction, the remainder of the content was written primarily by Claude-Opus-5.5 (medium effort, in Claude Code), with heavy editorial input from myself and GPT-6-Astra (high effort, in Codex). The paper describes technical work produced and formalized primarily by the same two models, in a collaboration enabled by javascript tools for them to send messages to each other. If you are primarily interested in how LLM agents were used to produce, formalize, and explain the result, please see Section \ref{sec:llm-usage}. 

I called our main result the Beating Sampling Theorem, because it shows how ``white-box'' statistical methods --- or arguably, mechanistic interpretability techniques --- are much more efficient than naive statistical sampling methods for estimating the frequency of rare events in neural network behavior.  Notably, mechanistic analysis is applied not only to the ``agent'' networks, but to the world (or world model) from which their inputs are drawn.

I was inspired to pursue this question by the Alignment Research Center (ARC), and by my own nagging sense that the mechanistic interpretability of the world should also be represented in mathematical intepretability research, as well as the cost of sampling from or analyzing the world. Previous work by ARC in this area has represented and analyzed primarily the mechanisms and sampling costs of the agent, not the world.  I also wanted to produce concrete, non-aysmptotic estimates of the computational costs and savings enabled by mechanistic analysis, at network scales somewhat comparable to current frontier AI models.

I'm not sure if everyone will consider the methods here ``mechanistic'' enough to be deserving of the term ``mechanistic estimation''.  In case of objections, I would offer the term ``perspicacious estimation'', as the estimation method peers into the innards of the network, and makes use of that insight to drastically reduce estimation costs.

\section{Motivation}\label{sec:motivation}

The statistical behavior of a neural network over an input distribution is usually estimated by sampling: collect inputs, evaluate the network on the inputs, count. However, behavior estimation by analyzing the internals of a network can be preferable to black-box sampling, for a variety of reasons:
\begin{enumerate}
\item \textbf{(sample count)} low-probability catastrophic events may require large numbers of samples to rule out, yielding high costs for eval-based safety assessments;
\item \textbf{(cost per sample)} for certain high-risk applications, access to real-world inputs is difficult, yielding a prohibitively high cost per sample; and
\item \textbf{(evaluation risk)} even in a testing environment, running a network with powerful cyberspace capabilities can compromise its containment, as in the infamous OpenAI–Huggingface intrusion.
\end{enumerate}

By contrast, mechanistic interpretability methods aim to understand a network's behavior by analyzing its internal structure after training. Alignment problems may thus be addressed by looking for sources of malignment inside the network and ruling them out, or validating the presence of alignment mechanisms. If successful, mechanistic safety approaches could hence reduce all three of the costs above: large sampling counts, high costs per sample, and evaluation risks.

This note exhibits a process for estimating the failure rate of a random MLP, under a random failure metric and an arbitrary threshold.  It follows on the work of the Alignment Research Center (ARC) on their ``Matching Sampling Principle'' \citep{arc-2025-competing,arc-2026-wide-mlps} as discussed in the Related Work section at the end.  Unlike ARC's asymptotic width guarantees, our guarantees use explicit network dimensions, and we analyze a non-Gaussian input distribution generated by another, larger neural network rather than the Gaussian inputs of \citet{arc-2026-wide-mlps}.

% The size-comparison table, shown in Section 3 and again in Appendix A.
\newcommand{\sizestable}{%
\setlength{\tabcolsep}{2.4pt}%
\begin{tabular}{@{}lrrrrr@{}}
\toprule
Network & $n$ & $h$ & $\ell$ & $n'$ & Parameters \\
\midrule
Kimi K3 MLPs (routed experts)~\citep{kimi-2026-k3} & $3{,}584$ & $2{,}752{,}512$ & $1$ & $3{,}584$ & $2.8\cdot10^{12}$ \\
GLM-5.3 MLPs~\citep{zai-2026-glm53} & $6{,}144$ & $524{,}288$ & $1$ & $6{,}144$ & $7.53\cdot10^{11}$ \\
DeepSeek V4 Pro MLPs~\citep{deepseek-2026-v4} & $7{,}168$ & $1{,}179{,}648$ & $1$ & $7{,}168$ & $1.6\cdot10^{12}$ \\
Qwen3.8-2.4T-A95B MLPs~\citep{qwen-2026-38} & $8{,}192$ & $1{,}048{,}576$ & $1$ & $8{,}192$ & $2.4\cdot10^{12}$ \\
Muse Spark Open MLPs~\citep{meta-2026-muse} & $4{,}096$ & $1{,}048{,}576$ & $1$ & $4{,}096$ & $8.16\cdot10^{11}$ \\
\midrule
\textbf{\bst, Agent ($\boldsymbol{A}$)} & $\boldsymbol{8{,}000}$ & $\boldsymbol{64{,}000{,}000}$ & $\boldsymbol{2}$ & $\boldsymbol{8{,}000}$ & $\boldsymbol{4.1\cdot10^{15}}$ \\
\textbf{\bst, World ($\boldsymbol{W}$)} & $\boldsymbol{16{,}000}$ & $\boldsymbol{4.58\cdot10^{13}}$ & $\boldsymbol{2}$ & $\boldsymbol{8{,}000}$ & $\boldsymbol{2.1\cdot10^{27}}$ \\
\bottomrule
\end{tabular}}


\section{Setting and result}\label{sec:setting}

\paragraph{MLP notation.} For positive integers $n,h,\ell,n'$, let $\mathrm{MLP}(n,h,\ell,n')$ denote the set of multi-layer perceptrons (MLPs) with ReLU activation functions and no biases, with input dimension $n$, hidden width $h$, $\ell$ hidden layers and output dimension $n'$.

\paragraph{Network sizes.} Rather than working asymptotically, our Beating Sampling Theorem will apply to agent MLPs above a certain size.  Table~\ref{tab:sizes} compares this size with the MLP components of open-weight frontier LLMs. Appendix~\ref{app:sizes} gives more details. Amongst these real-world models, $h\approx n^{1.6}$. The Beating Sampling Theorem uses $h=n^2$.

\begin{table}[h]
\centering\small
\sizestable
\caption{The agent and world networks of the \bst{} at $n=8{,}000$, and the MLP
blocks of large open-weight transformers. For a mixture-of-experts layer,
$h$ is the total width of its routed experts. The last column counts the
parameters of the whole model; see Appendix~\ref{app:sizes}.}
\label{tab:sizes}
\end{table}

\paragraph{The agent and world networks.}  Consider an ``agent'' $A\in\mathrm{MLP}(n,n^2,2,n)$.  We also need some way of representing the complexity of the broader world from which the agent will receive inputs.  For that, consider a much larger ``world'' network $W\in\mathrm{MLP}(2n,\lceil n^{3.5}\rceil,2,n)$.  This world and agent will not interact sequentially, so the ``agent'' is not very agent-like; these terms were chosen mainly to represent their relative sizes.  Their dimensions were chosen somewhat arbitrarily to achieve about $n^3$ (roughly a trillion) times more parameters in the world network than in the agent network when we later prove the \bst{}.  The world network will transform a high-entropy input $u$ into a normalized input $W(u)/|W(u)|$ for the network $A$, where $u$ is sampled uniformly from the unit sphere of $\RR^{2n}$.  The normalization step reflects the practice of neural network developers choosing reasonably bounded units in which the networks receive their inputs from the world.

Now suppose the agent network $A$ is being evaluated by a smaller ``failure metric'' network $F\in\mathrm{MLP}(n,\lceil n^{1.5}\rceil,2,1)$, with a threshold $t\in\RR$ defining a Boolean failure classifier $C(y)\coloneqq \mathbf 1[F(y)>t]$. Suppose the threshold $t$ is arbitrary but fixed, and $(W,A,F)$ are drawn independently with independent centered Gaussian weights, where the variance of each weight is $2/d$, and $d$ is the number of inputs to the layer containing it (He initialization).

Given $(W,A,F,t)$, we wish to estimate the failure probability 
\[p\coloneqq \PP_u\left[C\left(A\left(\frac{W(u)}{|W(u)|}\right)\right)=1\right]\]
with $u$ sampled from the unit sphere as above. The Beating Sampling Theorem will give a concrete procedure $\Est(W,A,F,r;t)$, which will use the weights of $W$ and $A$ but never evaluate either network on an input; here $r$ is a random seed. Even accounting for $W$, $\Est(W,A,F,r;t)$ will use fewer than $1.0003\cdot2n^7$ arithmetic operations, which is about the cost of one exact sample of $C(A(x))$. For every $n\ge8{,}000$ and every $t$, for at least $99\%$ of inputs $(W,A,F,r)$, $\Est(W,A,F,r;t)$ will return an upper bound $U$ with $p\le U\le\max(10p,1/100)$. The success-probability and operation-count bounds of the \bst{} are proved in Lean, for the definitions and cost model used here.

\paragraph{Success and cost.}
Let $x_0\coloneqq(1,0,\dots,0)\in\RR^n$ be the first standard basis vector. When $|W(u)|>0$ let
$x\coloneqq W(u)/|W(u)|$, else let $x\coloneqq x_0$. Then $p=\PP_u[F(A(x))>t]$. A program
outputting $U$ \emph{succeeds} if $p\le U\le\max(10p,1/100)$~\leanmark{1}. Its \emph{cost} is the
number of operations it performs on real numbers: additions, subtractions,
multiplications, divisions, square roots, $\relu$s, comparisons, branches and
table reads at computed indices; here $v/0\coloneqq0$, $\sqrt v\coloneqq$ the nonnegative square
root of $\max(0,v)$, and $\relu(v)\coloneqq\max(0,v)$ for $v\in\RR$. Reading an input (an entry of the seed $r$ or a weight of $W$, $A$ or $F$), a stored intermediate value or a rational constant costs nothing in the cost model in Lean~\leanmark{2}, although future work could easily include these costs, as they are negligible (see the program below).
\begin{leanitems}
\leanitem[1]{Part157Defs,Part156Out}; \href{\leanrepo}{\leanstepA}.
\leanitem[2]{Part130Cost,Part282Join,Part288Search,Part289Run}.
\end{leanitems}

\phantomsection\label{thm:main}
\begin{bstheorem}
Let $n\ge 8{,}000$ be an integer and $t\in\RR$. Then $\Est(W,A,F,r;t)$ applies
neither $W$ nor $A$ to any input, performs
fewer than $1.0003\cdot 2n^7$ operations, and for at least $99\%$ of inputs
$(W,A,F,r)$ its output $U$ satisfies
\[
  p\;\le\;U\;\le\;\max(10p,\,1/100).
\]
\end{bstheorem}

For comparison, collecting a single $C(A(x))$ sample by evaluating $W$, $A$ and $F$ layer by layer needs at least $2n^7$
operations (Section~\ref{sec:program}). Both bounds of the \bst{} are checked in Lean, using only the axioms
\texttt{propext}, \texttt{Classical.choice} and \texttt{Quot.sound}~\leanmark{1}.
\begin{leanitems}
\leanitem[1]{Part293FinalAudit,Part289Cost}; declarations \texttt{Part293.classifier\_stated2\_8000} and \texttt{Part289.progOut\_run2}; \href{\leanrepo}{\nolinkurl{BSP}}.
\end{leanitems}

\section{Comparison with estimation by naive sampling}\label{sec:comparison}

The \bst{} says that $\Est$ costs less than $1.0003$ times as much as one
exact sample of $C(A(x))$. How many samples would we need if we estimated 
$p$ just by naively sampling and tallying $C$-values?

Drawing independent inputs $u$ and naively evaluating $C(A(x))$ on each gives independent samples, like coin flips, each equal to
$1$ with probability $p$. Recall that a program outputting $U$
\emph{succeeds} if $p\le U\le\max(10p,1/100)$ (Section~\ref{sec:setting}). We call a program a \emph{naive-99\% sampler} if it sees only these flips and
its own independent random numbers, decides after each flip whether to stop,
and stops with probability $1$ and succeeds with probability at least $0.99$,
for \emph{every} value $p\in[0,1]$. Such a program is naive in three ways: it
evaluates every network in full on every input; it ignores the
weights of $W$, $A$ and $F$, which it uses only through the flips; and it
insists on succeeding with probability $0.99$ for every $p$, instead
of for $99\%$ of inputs $(W,A,F,r)$ as $\Est$ will.

Let $N$ be the number of flips that a naive-99\% sampler uses. For
$s\in[0,1]$ and $v\in(0,1)$, let
$\mathrm{KL}(s\Vert v)\coloneqq s\ln(s/v)+(1-s)\ln((1-s)/(1-v))$, with
$0\ln0\coloneqq0$. Say that an output $U$ \emph{succeeds at} a flip
probability $s$ if $s\le U\le\max(10s,1/100)$.

\begin{lemma}\label{lem:wald}
Let $p\in[0,1]$ and $q\in(0,1)$, and suppose that no $U$
succeeds both at $p$ and at $q$. Then every naive-99\% sampler satisfies
$\EE_pN\ge0.98\ln99/\mathrm{KL}(p\Vert q)$.
\end{lemma}
\begin{proof}
Lemma~\ref{lem:wald} extends to randomized samplers Wald's lower bound on
the expected number of observations of a sequential
test~\citep[Section~4.7, eq.~(4.80)]{wald-1945}. In the present terms,
Wald's bound states the following. Let a non-randomized sequential test
between flip probabilities $p$ and $q$ stop with probability $1$ at both,
and let $a$ and $b$ be its error probabilities at $p$ and at $q$. Then
$\EE_pN\cdot\mathrm{KL}(p\Vert q)\ge\mathrm{KL}(1-a\Vert b)$, where
we extend $\mathrm{KL}$ to $v\in\{0,1\}$ by $s\ln(s/0)\coloneqq\infty$ for
$s>0$ (Wald does not treat zero error probabilities; the bound holds trivially there). (Wald writes $-\mathrm{KL}(p\Vert q)$ as $E_0(z)$.)

The outputs that succeed at
$p$ and the outputs that succeed at $q$ form two disjoint nonempty intervals,
so one interval lies above the other. Let the sampler decide for the flip
probability of the upper interval if $U$ lies above the lower interval, and
for the other flip probability otherwise. At each of $p$ and $q$, success
implies the correct decision, so both error probabilities $a$ and $b$ are at
most $0.01$. Let $V$ be the sampler's own random numbers, which are
independent of the flips. For almost every value $v$ of $V$, the sampler with
$V=v$ is a non-randomized sequential test that stops with probability $1$ at
both $p$ and $q$. Let $a(v)$ and $b(v)$ be its error probabilities, so that
$\EE a(V)=a$ and $\EE b(V)=b$. Wald's bound gives
$\EE_p[N\mid V=v]\cdot\mathrm{KL}(p\Vert q)\ge\mathrm{KL}(1-a(v)\Vert b(v))$.
$\mathrm{KL}$ is jointly convex in its two arguments, so averaging over $V$
gives $\EE_pN\cdot\mathrm{KL}(p\Vert q)\ge\mathrm{KL}(1-a\Vert b)$. For
$s\ge v$, $\mathrm{KL}(s\Vert v)$ increases in $s$ and decreases in $v$; since
$1-a\ge0.99\ge0.01\ge b$, this gives
$\mathrm{KL}(1-a\Vert b)\ge\mathrm{KL}(0.99\Vert0.01)=0.98\ln99$. Finally,
$p\ne q$, since otherwise $U=p$ would succeed at both, so
$\mathrm{KL}(p\Vert q)>0$. Hence
$\EE_pN\ge0.98\ln99/\mathrm{KL}(p\Vert q)$, as claimed.
\end{proof}

For $p\le1/1000$, Lemma \ref{lem:wald} implies that every naive-99\% sampler needs more than $448$ samples on
average, as follows. In this range, $\max(10p,1/100)=1/100$, so success at
$p$ requires $U\le1/100$. For every $q\in(1/100,1)$, success at $q$ requires
$U\ge q>1/100$. So Lemma~\ref{lem:wald} gives
$\EE_pN\ge0.98\ln99/\mathrm{KL}(p\Vert q)$, and letting $q$ decrease to $1/100$
gives $\EE_pN\ge0.98\ln99/\mathrm{KL}(p\Vert1/100)$. The denominator decreases
in $p$ on $[0,1/100)$, so the bound is least at $p=0$, where it equals
$0.98\ln99/\ln(100/99)>448$. Hence, for $p\le1/1000$, every naive-99\%
sampler needs more than $448$ samples on average, as claimed.
Each flip is a naively computed exact sample of $C(A(x))$, which costs at
least $2n^7$ operations~\leanmark{1}. So when $p\le1/1000$, every naive-99\%
sampler has expected cost more than $448\cdot2n^7$, which is more than $447$
times the cost bound $1.0003\cdot2n^7$ of $\Est$.

\begin{leanitems}
\leanitem[1]{Part189Baseline}. The bounds from Wald's inequality and
their numerical values were not checked in Lean.
\end{leanitems}


Now, for $1/1000<p<1/10$, every naive-99\% sampler needs more than $31$ samples on
average, and for $p\ge1/10$ more than $1.95$ samples, as follows. If
$1/1000<p<1/10$, success at $p$ requires $U\le10p$, while success at
$q$ requires $U\ge q>10p$ for every $q\in(10p,1)$. If $p>1/100$, success at
$p$ requires $U\ge p$, while success at $q$ requires
$U\le\max(10q,1/100)<p$ for every $q\in(0,p/10)$. So Lemma~\ref{lem:wald} applies to these pairs.
Letting $q$ decrease to $10p$ or increase to $p/10$ in Lemma~\ref{lem:wald}
gives $\EE_pN\ge0.98\ln99/\mathrm{KL}(p\Vert10p)$ or
$\EE_pN\ge0.98\ln99/\mathrm{KL}(p\Vert p/10)$, respectively. Both bounds
decrease in $p$. Hence $\EE_pN>31.16$ for $1/1000<p<1/10$ and $\EE_pN>1.955$
for $p\ge1/10$, using the first bound up to $p=1/100$ and the second above it,
as claimed.
Since each flip costs at least $2n^7$ operations, the expected cost of every
naive-99\% sampler
is more than $31$ times the cost bound of $\Est$ when $1/1000<p<1/10$, and
more than $1.95$ times it when $p\ge1/10$. 


In summary, every naive-99\% sampler has a higher expected cost than $\Est$
at every $p$: more than $1.95$ times for $p\ge1/10$, more than $31$ times for
$1/1000<p<1/10$, and more than $447$ times for $p\le1/1000$. This trend underscores the
importance of using statistical tests that are ``perspicacious'' in that they
can read and use network weights, and/or ``pragmatic'' in that they work on a large
fraction of inputs instead of all inputs, especially for low-probability
high-cost events.

\section{The program $\Est$ and the proof of the \bst{}}\label{sec:program}

$\Est$ takes the input $(W,A,F,r;t)$ and follows the order of the networks'
own computation, i.e., through $W$, then $A$, then $F$, with draws from the seed $r$ in place of $u$~\leanmark{1}.
Each step below says what $\Est$ computes, then analyzes it, stating each
lemma where the analysis establishes it, and then says what the step
costs~\leanmark{2}. The closing paragraphs combine the lemmas into the \bst{}.
\begin{leanitems}
\leanitem[1]{Part157Paper,Part169Pass,Part272Stated,Part289Design};
  \href{\leanrepo}{\leanstepB}.
\leanitem[2]{Part282World,Part282Agent2,Part150Prog4,Part282RowData,Part282Grid,Part282NodeFull,Part282NodeAny,Part282Draw,Part282Out,Part282JoinCost,Part289Hit,Part289Run,Part289Cost}.
\end{leanitems}

\paragraph{Notation.} For vectors, $|\cdot|$ is the Euclidean norm, $\langle\cdot,\cdot\rangle$
the inner product, and $\cos\angle(y,z)\coloneqq \langle y,z\rangle/(|y||z|)$. $I$ is
the identity matrix of the size the context requires. $\SSS^{k-1}$ is the
unit sphere of $\RR^k$. A subscript on $\PP$ or
$\EE$ names the variables averaged over, with the others fixed. An
unsubscripted $\PP$ or $\EE$ is over $(W,A,F,r)$.

\paragraph{Weight matrices.}
With $\relu$ applied coordinatewise to real vectors, for $z\in\RR^{2n}$ and $z'\in\RR^n$ write
\[W(z)\coloneqq W_3\relu(W_2\relu(W_1z)),\qquad A(z')\coloneqq A_3\relu(A_2\relu(A_1z')),\]
and $m_W\coloneqq \lceil n^{3.5}\rceil$~\leanmark{1}.
\begin{leanitems}
\leanitem[1]{Part130Defs,Part157Triple}.
\end{leanitems}

%AC: reviewed up to here.

\subsection{High-level overview}
\begin{lemma}\label{lem:samplecost}
One exact sample of $C(A(x))$, computed by evaluating $W$, $A$ and $F$ layer by
layer, costs at least $2n^7$ operations, which is more
than $4.1\cdot10^{27}$ at $n=8{,}000$.
\end{lemma}
\begin{proof}
An exact sample computes the product of $W_2$ and the world's first hidden
vector at $u$. Since $W_2$ is an $m_W\times m_W$ matrix, this product takes
$2m_W^2-m_W$ operations. The product $W_1u$ takes more than $m_W$ further
operations. Hence an exact sample costs at least $2m_W^2\ge2n^7$
operations, since $m_W\ge n^{3.5}$, as claimed.
\end{proof}
$\Est$ multiplies $W_2$ by one vector only: an estimate of the world's mean first
hidden vector. $\Est$ never applies $W$ or $A$ to an input.

\begin{lemma}\label{lem:gauss}
Let $M$ be a random matrix with independent $\Nn(0,\sigma^2)$ entries, and let
$D$ be random data independent of $M$. Let $g$ be a fixed measurable function
from the values of $D$ to vectors, and let $v\coloneqq g(D)$. Let $\Pi_v$ be the orthogonal
projection onto the directions orthogonal to $v$, i.e.,
$\Pi_v\coloneqq I-vv^\top/|v|^2$ if $v\ne0$, and $\Pi_0\coloneqq I$. Then, given
$D$, the matrix $M\Pi_v$ is Gaussian and independent of $Mv$~\leanmark{1}.
\end{lemma}
\begin{proof}
Lemma~\ref{lem:gauss} is standard; see \citet[arXiv v3, Lemmas~G.6--G.7]{yang-2019-tp1},
which extend \citet[arXiv v4, Lemmas~11--12]{bayati-montanari-2011} to singular constraints. First let $v$ be fixed. Each entry of $M\Pi_v$ and of $Mv$ is a
linear combination of the entries of $M$, so the pair $(M\Pi_v,Mv)$ is
jointly Gaussian with mean zero. The rows of $M$ are independent
$\Nn(0,\sigma^2I)$ vectors, so the cross-covariances are
$\EE[(M\Pi_v)_{ij}(Mv)_k]=\delta_{ik}\sigma^2(\Pi_vv)_j=0$, since
$\Pi_vv=0$. Jointly Gaussian vectors with zero cross-covariance are
independent, so $M\Pi_v$ is independent of $Mv$. Now let $D$ be random.
Since $D$ is independent of $M$, the distribution of $M$ given $D$ is the
distribution of $M$, and $v$ is fixed given $D$. Hence, by the case of a fixed $v$ applied to
each value of $D$, given $D$ the matrix $M\Pi_v$ is Gaussian and
independent of $Mv$, as claimed.
\end{proof}
\begin{leanitems}
\leanitem[1]{Part134Split,Part157Reg}.
\end{leanitems}

Lemma~\ref{lem:gauss} applied first to the world and then to the agent
implies that $\Est$ can avoid fully evaluating $W$ or $A$ on any inputs.

Applied to the world, Lemma~\ref{lem:gauss} gives a proxy input, as follows. The only use $\Est$ makes
of the last layer $W_3$ is to compute one vector $\chat\in\RR^n$. Step~2 will
compute $\chat\coloneqq W_3b_W$ for an estimate $b_W$ of the world's mean
last hidden vector. Call $\Omega_W\coloneqq W_3\Pi_{b_W}$ the residual part of
$W_3$. By Lemma~\ref{lem:gauss}, $\Omega_W$ is
Gaussian, even given $W_1$, $W_2$ and $\chat$, which are all that $\Est$ uses of
$W$. So for each $u$, over
$\Omega_W$, the input $x$ is the direction of the sum of a multiple of
$\chat$ and isotropic Gaussian noise. The scale of this noise depends on
$u$. $\Est$ replaces this scale by an estimate $\tau_0$ that is the same for every
$u$. Hence, with this replacement,
Lemma~\ref{lem:gauss} applied to the world gives a \emph{proxy input} $\tilde x$, whose
distribution depends only on $\chat$ and $\tau_0$, which $\Est$ computes, as claimed.
Lemma~\ref{lem:noisy} in step~2 will state the distribution of $x$ precisely.

Applied to the agent, Lemma~\ref{lem:gauss} gives the distribution of the agent's output,
as follows.
The only use $\Est$ makes of $A_3$ is to compute $c_A\coloneqq A_3b_A$ for an
estimate $b_A$ of the agent's mean last hidden vector (step~3). Call $\Omega_A\coloneqq A_3\Pi_{b_A}$ the residual part of $A_3$.
By Lemma~\ref{lem:gauss}, $\Omega_A$ is Gaussian, even given $W$, $A_1$, $A_2$ and
$c_A$. Hence, over $\Omega_A$, the agent's output at each input is the sum
of a multiple of $c_A$ and isotropic Gaussian noise, as claimed.
Lemma~\ref{lem:agentout} in step~3 will state the distribution of the agent's output precisely.

$\Est$ computes $b_W$ and $b_A$ from the earlier layers alone, so Lemma~\ref{lem:gauss}
applies to $b_W$ and $b_A$. The proof will not assume that $b_W$ and $b_A$ are accurate.
Their accuracy will enter only
through the error bounds $\Delta_W$ and $\Delta_A$ of the chain~\eqref{eq:chain} below and through the frame bounds of Lemma~\ref{lem:frame}.

At the proxy input, the coefficient of $c_A$ and the variance of the noise
in the agent's output of Lemma~\ref{lem:agentout} depend mainly on one number: the component of $\tilde x$
along $\chat/|\chat|$, which step~4 will call the \emph{height} of $\tilde x$. So
steps~4--5 will tabulate the coefficient and the variance at $512n$ heights
that depend only on the data $\Dd$ defined below, called \emph{nodes}. Step~6 will then compute $K=10^{12}$
\emph{proxy outputs} from the tables and the seed $r$. $\Est$
evaluates $F$ once on each proxy output, instead of evaluating $W$, $A$ and
$F$ once per exact sample. Finally, step~7 will scale and clip the fraction $R$
of proxy outputs that fail, and output the result.

Let $\Dd\coloneqq (W_1,W_2,\chat,A_1,A_2,c_A,F)$ be the data that $\Est$
uses apart from the seed $r$; $\Est$ reads $W_3$ and $A_3$ only to compute $\chat$ and $c_A$. Let $\tilde p$, which step~3 will define, be the failure probability at
the proxy input, and let $\tilde q$, which step~6 will define, be the
probability that a proxy output fails. The proof follows a chain of approximations from $p$
to $R$:
\begin{equation}\label{eq:chain}p\ \underset{\text{conc.}}{\approx}\ \EE_{\Omega_W}p\ \underset{\Delta_W}{\approx}\ \tilde p
\ \underset{\text{conc.}}{\approx}\ \EE_{\Omega_A}\tilde p\ \underset{\Delta_A}{\approx}\ \tilde q
\ \underset{\Delta_s}{\approx}\ R.\end{equation}
The two links marked ``conc.''\ will be concentration bounds over
$\Omega_W$ and $\Omega_A$. The other three will be bounds on the
\emph{square-root scale}, i.e., of the form $|\sqrt a-\sqrt b|\le\Delta$, with errors $\Delta_W$, $\Delta_A$ and $\Delta_s$. Step~6 will prove the
$\Delta_A$ link, and step~7 the $\Delta_s$ link. The paragraph ``Bounds over
the residual parts'' will prove the $\Delta_W$ link and the concentration bounds for the two
links marked ``conc.'', on a good set $\Gg$ of values of $\Dd$. The paragraph ``Assembly''
will combine all five links.

\subsection{Detailed overview}
\begin{enumerate}[leftmargin=*,itemsep=8pt,topsep=10pt plus 2pt]
\item \textbf{Constants for the layers of $W$.} This step computes
  constants that depend only on $n$.

  \emph{Computation:} Let $N\coloneqq 2n$. For any integer $T\ge1$ and
  $v\in\RR$, let $P_T$ and $Q_T$ be the truncated Taylor series of $\arctan$
  and $\arcsin$, i.e.,
  \[P_T(v)\coloneqq \sum_{0\le j<T}\frac{(-1)^jv^{2j+1}}{2j+1},\qquad
  Q_T(v)\coloneqq \sum_{0\le j<T}\binom{2j}{j}\frac{v^{2j+1}}{4^j(2j+1)}.\]
  Let $T_\pi$ be the least integer $T\ge1$ with $25^T\ge10^{23}N^2$.
  Compute Machin's approximation of $\pi$, i.e.,
  $\tilde\pi\coloneqq 16P_{T_\pi}(1/5)-4P_{T_\pi}(1/239)$. Compute an
  approximation of the mean of $\relu(u_1)$, i.e.,
  $\tilde c\coloneqq \tilde\pi^{-1}\prod_{j=1}^{n-1}\frac{2j}{2j+1}$. Compute
  an approximation of twice the variance of $\relu(u_1)$, i.e.,
  $\tilde d\coloneqq N^{-1}-2\tilde c^2$. Compute an approximation of the correlation of the
  first-layer features of two independent inputs, i.e.,
  $\tilde r\coloneqq 2N\tilde c^2$. Let $T_\kappa$ be the least integer
  $T\ge1$ with $0.355\cdot0.1014^T\le10^{-5}/n^2$. Compute an approximation
  of the first-order arc-cosine kernel at $\tilde r$, i.e.,
  \[\tilde\kappa\coloneqq \bigl(\sqrt{1-\tilde r^2}+\tilde r(\tilde\pi/2+Q_{T_\kappa}(\tilde r))\bigr)/\tilde\pi.\]

  \emph{Analysis:} The constant $\tilde\pi$ truncates the series of Machin's
  formula $\pi=16\arctan(1/5)-4\arctan(1/239)$, and the error $|\tilde\pi-\pi|$ falls
  geometrically in $T_\pi$.
  \begin{lemma}\label{lem:pi}
    $|\tilde\pi-\pi|\le4\cdot25^{-T_\pi}$~\leanmark{1}.
  \end{lemma}
  \begin{proof}
  Lemma~\ref{lem:pi} is standard: its proof combines Machin's formula with a
  bound on the tail of the Taylor series of $\arctan$. Let
  $T\coloneqq T_\pi$ and $0\le v<1$. Since
  $\arctan v=\sum_{j\ge0}(-1)^jv^{2j+1}/(2j+1)$, the tail after $T$ terms is
  at most the geometric sum $\sum_{j\ge T}v^{2j+1}$ in absolute value, i.e.,
  $|\arctan v-P_T(v)|\le v^{2T+1}/(1-v^2)$. By Machin's formula,
  $\tilde\pi-\pi=16(P_T(1/5)-\arctan\tfrac15)-4(P_T(1/239)-\arctan\tfrac1{239})$.
  Hence, by this tail bound at $v=1/5$ and at $v=1/239$, and since
  $239^{-2T}\le25^{-T}$,
  \[|\tilde\pi-\pi|\le\frac{16/5}{1-1/25}\,25^{-T}+\frac{4/239}{1-1/239^2}\,25^{-T}
  =\Bigl(\frac{10}{3}+\frac{956}{57{,}120}\Bigr)25^{-T}<4\cdot25^{-T},\]
  as claimed.
  \end{proof}
  \begin{leanitems}
  \leanitem[1]{Part157Paper,Part159Pi}.
  \end{leanitems}
  The constants $\tilde c$ and $\tilde d$ approximate $\EE_u\relu(u_1)$ and
  $2\operatorname{Var}_u\relu(u_1)$, respectively.
  \begin{lemma}\label{lem:cd}
    Let $c\coloneqq\pi^{-1}\prod_{j=1}^{n-1}\frac{2j}{2j+1}$ and
    $d\coloneqq N^{-1}-2c^2$, i.e., $\tilde c$ and $\tilde d$
    with $\pi$ in place of $\tilde\pi$. Then $c=\EE_u\relu(u_1)$ and
    $d=2\operatorname{Var}_u\relu(u_1)$~\leanmark{1}.
  \end{lemma}
  \begin{proof}
  Lemma~\ref{lem:cd} is an elementary computation. Since $u$ is uniform on $\SSS^{N-1}$, the coordinate $u_1$ has density
  proportional to $(1-s^2)^{(N-3)/2}$ on $(-1,1)$. Substituting
  $s=\sin\theta$, and using $N-2=2n-2$, gives
  \[\EE_u\relu(u_1)=\frac{\int_0^{\pi/2}\sin\theta\cos^{2n-2}\theta\,d\theta}
  {2\int_0^{\pi/2}\cos^{2n-2}\theta\,d\theta}
  =\frac{1/(2n-1)}{\pi\prod_{j=1}^{n-1}\frac{2j-1}{2j}},\]
  by Wallis' formula $\int_0^{\pi/2}\cos^{2n-2}\theta\,d\theta
  =\frac\pi2\prod_{j=1}^{n-1}\frac{2j-1}{2j}$, which follows from
  \citet[eqs.~5.12.1--5.12.2]{dlmf}, with $\Gamma(1/2)=\sqrt\pi$ and $\Gamma(v+1)=v\Gamma(v)$. Since
  $(2n-1)\prod_{j=1}^{n-1}\frac{2j-1}{2j+1}=1$, the right side of the display equals $c$, so
  $\EE_u\relu(u_1)=c$.
  Next, $\relu(s)^2+\relu(-s)^2=s^2$ for $s\in\RR$, the map $u\mapsto-u$
  preserves the distribution of $u$, and $\EE_uu_1^2=1/N$, since
  $\sum_iu_i^2=1$ and the coordinates of $u$ have the same distribution. So
  $\EE_u\relu(u_1)^2=1/(2N)$. Hence
  $2\operatorname{Var}_u\relu(u_1)=N^{-1}-2c^2=d$, as claimed.
  \end{proof}
  \begin{leanitems}
  \leanitem[1]{Part159CN,Part137Hat2}. Lean checks the formula for $c$ and
  the second moment $\EE_u\relu(u_1)^2=1/(2N)$; the identity for $d$ then
  follows by the last line of the proof.
  \end{leanitems}
  The constant $\tilde\kappa$ approximates the first-order arc-cosine kernel of
  \citet[eqs.~(3) and~(6)]{cho-saul-2009} at unit vectors with inner
  product $v$, i.e., $\kappa(v)\coloneqq (\sqrt{1-v^2}+v(\pi/2+\arcsin
  v))/\pi$ for $v\in[-1,1]$, at $\tilde r$. $\tilde\kappa$ replaces $\pi$ and $\arcsin$ by
  $\tilde\pi$ and $Q_{T_\kappa}$, with an error that shrinks as $T_\kappa$ grows.

  \emph{Cost:} none for $\tilde c$, $\tilde d$ and $\tilde\pi$ in the cost
  model in Lean, which charges nothing for rational constants; the square
  root, addition and division that form $\tilde\kappa$ are charged in step~2. For $8{,}000\le n\le16{,}384$, $T_\pi=23$ and
  $T_\kappa\le14$. Computing $\tilde c$, $\tilde\pi$ and $\tilde\kappa$ would
  then take at most $4n+250$ operations: about $4n$ for the product in
  $\tilde c$, about $6T_\pi\le150$ for $\tilde\pi$, and about
  $6T_\kappa\le90$ for $\tilde\kappa$. So $4n+250\le65{,}786$. Since
  $n\ge8{,}000$, also $4n+250<10^{-23}\cdot2n^7$.

\item \textbf{The proxy input.} This step computes the vector $\chat$ and
  the proxy noise scale $\tau_0$, which together define the proxy input. This step is the only
  one that multiplies by $W_2$.

  \emph{Computation:} Let $v_1,\dots,v_{m_W}\in\RR^N$ be the rows of $W_1$.
  Compute the estimated mean of each first-layer unit, i.e.,
  $\psi_i\coloneqq \tilde c|v_i|$ for $1\le i\le m_W$. Compute the mean
  squared row norm, i.e., $\Sigma_W\coloneqq m_W^{-1}\sum_i|v_i|^2$. Use $\Sigma_W$ to
  estimate the spread of the second-layer pre-activations, i.e.,
  $\varsigma_W\coloneqq \sqrt{\Sigma_W\tilde d}$. For any integer $J\ge1$ and
  any function $\Phi:\RR\to\RR$, let $B_J[\Phi]$ be the mean of $\Phi$ at a
  sum of $J$ independent uniform random signs divided by $\sqrt J$, i.e.,
  \[B_J[\Phi]\coloneqq \sum_{j=0}^J2^{-J}\binom Jj\,\Phi\bigl((2j-J)/\sqrt J\bigr).\]
  Let $J_W\coloneqq \lceil\sqrt{m_W}\rceil$, and let $\nu_W\coloneqq 2/m_W$ be
  the variance of an entry of $W_2$ and $W_3$. Compute the estimated mean
  $b_W$ of the world's last hidden layer, the image $\chat$ of $b_W$ under $W_3$, the estimated
  cosine $\ell_0$ between a last hidden vector of the world and $b_W$, and
  the proxy noise scale $\tau_0$, i.e.,
  \begin{gather*}
  b_{W,i}\coloneqq B_{J_W}[z\mapsto\relu((W_2\psi)_i+\varsigma_Wz)]\ \text{for }1\le i\le m_W,\qquad
  \chat\coloneqq W_3b_W,\\
  \ell_0\coloneqq \sqrt{\frac{\Sigma_W\tilde\kappa^2}{N\nu_W|b_W|^2}},\qquad
  \tau_0\coloneqq \sqrt{\nu_W|b_W|^2\max\bigl(0,\,\tfrac1{\ell_0^2}-1\bigr)}.
  \end{gather*}

  \emph{Analysis:} Let $h_W:\RR^N\to\RR^{m_W}$ be the map to the world's
  last hidden layer, i.e., $h_W\coloneqq \relu\circ W_2\circ\relu\circ W_1$.
  The estimates are computed layer by layer. The first-layer unit $i$ has mean
  $\EE_u\relu(v_i\cdot u)=|v_i|\,\EE_u\relu(u_1)$, and $\psi_i$ estimates this mean.
  The vector $b_W$ estimates $\EE_uh_W(u)$ by treating each second-layer
  pre-activation as a Gaussian with mean $(W_2\psi)_i$ and variance
  $\varsigma_W^2$. The rule $B_J$ approximates each such Gaussian expectation:
  $B_J[\Phi]\to\EE_Z\Phi(Z)$ as $J\to\infty$, for $Z\sim\Nn(0,1)$ and
  continuous $\Phi$ of polynomial growth. The arc-cosine kernel motivates
  $\ell_0^2$ as an estimate of $\cos^2\angle(h_W(u),b_W)$ for most $u$. The
  proof will use about $\psi$, $b_W$ and $\ell_0$ only the
  bound $\Delta_W$, the frame bound of Lemma~\ref{lem:frame}, and the events
  that define the good set $\Gg$ of ``Bounds over the residual parts''.

  Let $\Omega_W$ be the residual part of $W_3$,
  i.e., $\Omega_W\coloneqq W_3\Pi_{b_W}$, with $\Pi_{b_W}$ as in
  Lemma~\ref{lem:gauss}. For each $u$, split
  $h_W(u)$ along $b_W$, i.e., $h_W(u)=\lambda_W(u)b_W+h_W^\perp(u)$ with
  $\lambda_W(u)\coloneqq \langle h_W(u),b_W\rangle/|b_W|^2$ and
  $h_W^\perp(u)\coloneqq h_W(u)-\lambda_W(u)b_W$.
  \begin{lemma}\label{lem:wsplit}
    For every $u$, $W(u)=\lambda_W(u)\chat+\Omega_Wh_W^\perp(u)$.
  \end{lemma}
  \begin{proof}
  Since $h_W^\perp(u)$ is orthogonal to $b_W$,
  $W_3h_W^\perp(u)=\Omega_Wh_W^\perp(u)$. Hence, since $W_3b_W=\chat$,
  $W(u)=W_3h_W(u)=\lambda_W(u)\chat+\Omega_Wh_W^\perp(u)$, as claimed.
  \end{proof}
  Lemma~\ref{lem:wsplit} gives the distribution of $x$. For each $u$ with
  $\lambda_W(u)>0$, let the \emph{world noise scale} be
  \[\tau_W(u)\coloneqq \frac{\sqrt{\nu_W}\,|h_W^\perp(u)|}{\lambda_W(u)}.\]
  \begin{lemma}\label{lem:noisy}
    For each $u$ with $\lambda_W(u)>0$, $x$ is the direction of a noisy copy
    of $\chat$, i.e., given $W_1$, $W_2$ and $\chat$, the distribution of $x$
    over $\Omega_W$ is that of $(\chat+\tau_W(u)\xi)/|\chat+\tau_W(u)\xi|$,
    with $\xi\sim\Nn(0,I)$ in $\RR^n$, and with $x_0$ in place of this
    quotient when $\chat+\tau_W(u)\xi=0$~\leanmark{1}.
  \end{lemma}
  \begin{proof}
  Fix $u$ with
  $\lambda_W(u)>0$. The vector $b_W$ is a fixed function of $(W_1,W_2)$, and
  $W_3$ is independent of $(W_1,W_2)$. So, by Lemma~\ref{lem:gauss} with
  $M=W_3$ and $D=(W_1,W_2)$, given $W_1$ and $W_2$, the matrix
  $\Omega_W=W_3\Pi_{b_W}$ is Gaussian and independent of $\chat=W_3b_W$. The
  rows of $\Omega_W$ are independent $\Nn(0,\nu_W\Pi_{b_W})$ vectors, and
  $\Pi_{b_W}h_W^\perp(u)=h_W^\perp(u)$, where $h_W^\perp(u)$ is a fixed function of
  $(W_1,W_2)$. So given $W_1$, $W_2$ and $\chat$, the second term of
  Lemma~\ref{lem:wsplit}, $\Omega_Wh_W^\perp(u)$, has the distribution
  $\Nn(0,\nu_W|h_W^\perp(u)|^2I)$, i.e., the distribution of
  $\sqrt{\nu_W}\,|h_W^\perp(u)|\,\xi$. Hence $W(u)$ has the distribution of
  $\lambda_W(u)(\chat+\tau_W(u)\xi)$. Multiplication by $\lambda_W(u)>0$
  changes neither the direction of a vector nor whether the vector is $0$.
  Hence $x$ has the distribution of
  $(\chat+\tau_W(u)\xi)/|\chat+\tau_W(u)\xi|$, with $x_0$ in place of this
  quotient when $\chat+\tau_W(u)\xi=0$, as claimed.
  \end{proof}
  \begin{leanitems}
  \leanitem[1]{Part157Reg,Part157Gp,Part157Repl,Part132Cos}. Lean takes
  $x_0=(1,0,\dots,0)$, and states the lemma as the law of the direction of
  $\lambda_W(u)\chat+\sigma_W(u)\gamma$, with
  $\sigma_W(u)\coloneqq\sqrt{\nu_W}\,|h_W^\perp(u)|$ and
  $\gamma\sim\Nn(0,I)$, and also averaged over $u$; that form needs no condition on
  $\lambda_W(u)$, and dividing by $\lambda_W(u)>0$ gives the form of the lemma.
  \end{leanitems}
  The world noise scale $\tau_W(u)$ depends on $u$ only through a cosine, i.e.,
  \[\tau_W(u)^2=\nu_W|b_W|^2\bigl(\cos^{-2}\angle(h_W(u),b_W)-1\bigr).\]
  Replacing the squared cosine by $\ell_0^2$ gives $\tau_0^2$, which $\Est$
  computes. Replacing $\tau_W(u)$ by $\tau_0$ in Lemma~\ref{lem:noisy} gives
  the \emph{proxy input}, i.e.,
  $\tilde x\coloneqq (\chat+\tau_0\xi)/|\chat+\tau_0\xi|$ if
  $\chat+\tau_0\xi\ne0$ and $\tilde x\coloneqq x_0$ otherwise, with
  $\xi\sim\Nn(0,I)$ in $\RR^n$. Unlike the distribution of $x$, the
  distribution of $\tilde x$ does not depend on the residual part $\Omega_W$: it
  depends only on $\chat$ and $\tau_0$, which $\Est$ computes. So $\Est$ can
  estimate averages over $\tilde x$ (step~5) and draw exact samples of its
  height (step~6).
  The bound $\Delta_W$ will cover the error of replacing $x$ by $\tilde x$, and also the
  rare inputs with $\lambda_W(u)\le0$.

  \emph{Cost:} $2m_W^2+4m_WJ_W+O(nm_W)$ operations. Almost all of these
  operations compute the product $W_2\psi$, which costs $2m_W^2-m_W$. At $n=8{,}000$,
  $J_W=6{,}767{,}177$ and $4m_WJ_W/(2m_W^2)<3\cdot10^{-7}$. The cost model
  in Lean does not charge for the weights $2^{-J_W}\binom{J_W}j$ of $B_{J_W}$.
  These weights $a_j$ satisfy $a_0=2^{-J_W}$ and $a_{j+1}=a_j(J_W-j)/(j+1)$.
  Computing all these weights by this recursion would take about $3J_W$ operations, i.e., about $2\cdot10^7$ at
  $n=8{,}000$ and $7\cdot10^7$ at $n=16{,}384$, less than $10^{-20}$ times
  $2m_W^2$.

\item \textbf{The agent's output distribution.} This step computes the vector
  $b_A$ along which $\Est$ splits the agent's last hidden layer, and the image
  $c_A$ of $b_A$ under $A_3$, as step~2 did for the world.

  \emph{Computation:} Let $a_1,\dots,a_{n^2}\in\RR^n$ be the rows of $A_1$, and
  let $\nu_A\coloneqq 2/n^2$ be the variance of an entry of $A_2$ and $A_3$.
  Compute the root-mean-square norm of the unnormalized proxy
  input, i.e., $\rho\coloneqq \sqrt{|\chat|^2+n\tau_0^2}$. For
  $1\le i\le n^2$, let $\phi_i$ be the activation of first-layer unit $i$ at the
  unnormalized proxy input, i.e., $\phi_i(z)\coloneqq \relu(a_i\cdot\chat+\tau_0|a_i|z)$
  for $z\in\RR$. Estimate the first two moments of $\phi_i$ at a standard
  Gaussian by
  $\hat h_i\coloneqq B_n[\phi_i]$ and $\hat h^{(2)}_i\coloneqq B_n[\phi_i^2]$.
  Estimate the spread of the second-layer pre-activations, i.e.,
  $\hat\varsigma\coloneqq \sqrt{\nu_A(\sum_i\hat h^{(2)}_i-|\hat h|^2)}$.
  Compute the estimated mean $b_A$ of the agent's last hidden layer and the
  image $c_A$ of $b_A$ under $A_3$, i.e.,
  \[b_{A,j}\coloneqq \rho^{-1}B_n[z\mapsto\relu((A_2\hat h)_j+\hat\varsigma z)]\ \text{for }1\le j\le n^2,\qquad
  c_A\coloneqq A_3b_A.\]

  \emph{Analysis:} Let $h:\RR^n\to\RR^{n^2}$ be the map to the agent's last
  hidden layer, i.e., $h\coloneqq \relu\circ A_2\circ\relu\circ A_1$. Since
  $A$ has no biases, $h(\tilde x)=h(\chat+\tau_0\xi)/|\chat+\tau_0\xi|$ whenever
  $\chat+\tau_0\xi\ne0$, and
  $\rho^2=\EE_\xi|\chat+\tau_0\xi|^2$. Each first-layer pre-activation
  $a_i\cdot(\chat+\tau_0\xi)$ has the distribution of
  $a_i\cdot\chat+\tau_0|a_i|Z$ with $Z\sim\Nn(0,1)$, and the second layer is
  as in step~2. So $b_A$ is designed to approximate $\EE_{\tilde x}h(\tilde x)$.

  As for the world in step~2, split the agent's last hidden layer along
  $b_A$. For $z\in\RR^n$, let
  $\lambda(z)$ be the coefficient of $h(z)$ along $b_A$, $h^\perp(z)$ the rest
  of $h(z)$, and $S(z)$ the noise variance, i.e.,
  \[\lambda(z)\coloneqq \frac{\langle h(z),b_A\rangle}{|b_A|^2},\qquad
  h^\perp(z)\coloneqq h(z)-\lambda(z)b_A,\qquad
  S(z)\coloneqq \nu_A|h^\perp(z)|^2.\]
  Let $\Omega_A$ be the residual part of $A_3$,
  i.e., $\Omega_A\coloneqq A_3\Pi_{b_A}$, with $\Pi_{b_A}$ as in
  Lemma~\ref{lem:gauss}.
  \begin{lemma}\label{lem:asplit}
    For every $z\in\RR^n$, $A(z)=\lambda(z)c_A+\Omega_Ah^\perp(z)$.
  \end{lemma}
  \begin{proof}
  Since $h^\perp(z)$ is orthogonal to $b_A$,
  $A_3h^\perp(z)=\Omega_Ah^\perp(z)$. Hence, since $A_3b_A=c_A$,
  $A(z)=A_3h(z)=\lambda(z)c_A+\Omega_Ah^\perp(z)$, as claimed.
  \end{proof}
  \begin{lemma}\label{lem:agentout}
    The agent's output is $c_A$ times a coefficient plus isotropic
    noise, i.e., let $z$ be a fixed function of $W$, $A_1$, $A_2$, $c_A$ and
    $\xi$, with values in $\RR^n$. Then, given $W$, $A_1$, $A_2$, $c_A$ and
    $\xi$, the distribution of $A(z)$ over $\Omega_A$ is that of
    $\lambda(z)c_A+\sqrt{S(z)}\,\eta$, with $\eta\sim\Nn(0,I)$ in
    $\RR^n$~\leanmark{1}.
  \end{lemma}
  \begin{proof}
  Lemma~\ref{lem:agentout} is the Gaussian conditioning trick with the single
  constraint $c_A=A_3b_A$; see \citet[arXiv v3, Lemma~G.7]{yang-2019-tp1}. The vector $b_A$ is a fixed function of $W$, $A_1$ and $A_2$, and
  $A_3$ is independent of $(W,A_1,A_2,\xi)$. So, by Lemma~\ref{lem:gauss}
  with $M=A_3$ and $D=(W,A_1,A_2,\xi)$, given $W$, $A_1$, $A_2$ and $\xi$,
  the matrix $\Omega_A$ is Gaussian and independent of $c_A=A_3b_A$, and the
  rows of $\Omega_A$ are independent $\Nn(0,\nu_A\Pi_{b_A})$ vectors. So
  given $W$, $A_1$, $A_2$, $c_A$ and $\xi$, the vector $z$ is fixed,
  $\Omega_A$ has the same distribution, and, since
  $\Pi_{b_A}h^\perp(z)=h^\perp(z)$, the term $\Omega_Ah^\perp(z)$ has the
  distribution $\Nn(0,\nu_A|h^\perp(z)|^2I)=\Nn(0,S(z)I)$. Hence, by
  Lemma~\ref{lem:asplit}, $A(z)$ has the distribution
  of $\lambda(z)c_A+\sqrt{S(z)}\,\eta$ given $W$, $A_1$, $A_2$, $c_A$ and $\xi$, as claimed.
  \end{proof}
  \begin{leanitems}
  \leanitem[1]{Part157Reg,Part157Net,Part157Triple}. Lean proves a joint
  form (\texttt{triple\_disint}): $\Dd$ and the three networks
  have the law obtained by rebuilding $W_3$ and $A_3$ from $\Dd$
  and from Gaussian matrices $G_W$ and $G_A$, independent of each other and of
  $\Dd$, whose entries have variances $\nu_W$ and $\nu_A$,
  respectively; the residual parts are $G_W\Pi_{b_W}$ and $G_A\Pi_{b_A}$. With
  \texttt{mlp2\_reb}, which proves Lemma~\ref{lem:asplit} at every input, and
  \texttt{resid\_gauss}, this gives Lemma~\ref{lem:agentout} at each input.
  \end{leanitems}
  The splitting of $A_3$ in Lemma~\ref{lem:asplit} is the Gaussian
  conditioning trick of \citet{bolthausen-2014}, also used by
  \citet{bayati-montanari-2011} and by \citet{yang-2019-tp1} in Tensor Programs. Let $\tilde p$ be the
  failure probability at the proxy input, i.e.,
  $\tilde p\coloneqq \PP_\xi[F(A(\tilde x))>t]$, a function of $\Dd$ and
  $\Omega_A$.

  \emph{Cost:} $2n^4+20n^3+O(n^2)$ operations, including $2n^3$ for $c_A=A_3b_A$. Almost all of these
  operations compute the product $A_2\hat h$. The cost model in Lean does not
  charge for the weights of $B_n$. Computing these weights as in step~2 would
  take about $3n$ more operations.

\item \textbf{A grid of heights.} By Lemma~\ref{lem:agentout}, the agent's output at $z$ is
  determined in distribution by $\lambda(z)$ and $S(z)$. This step defines the
  \emph{nodes}: a grid of $512n$ heights. Step~5 will tabulate estimates of
  two hidden-layer averages at the nodes, and step~6 will use these estimates to estimate $\lambda$ and
  $S$.

  \emph{Computation:} Let $e_1$ be the first standard basis vector of $\RR^n$,
  and let $\hat u$ be the direction of $\chat$, i.e.,
  $\hat u\coloneqq \chat/|\chat|$ if $\chat\ne0$, else
  $\hat u\coloneqq -e_1$. Call $z\cdot\hat u$ the \emph{height} of a unit
  vector $z\in\RR^n$. Compute the central height $s_c\coloneqq |\chat|/\rho$
  and the \emph{nodes}, i.e., the $512n$ heights around $s_c$ given by
  $s_g\coloneqq s_c+(6\tau_0/\rho)\bigl((2g+1)/(512n)-1\bigr)$ for integers
  $0\le g<512n$. At $n=8{,}000$, there are $512n=4{,}096{,}000$ nodes. Compute the
  component of each row $a_i$ along $\hat u$, and the length of the component of
  $a_i$ orthogonal to $\hat u$, i.e.,
  $\alpha_i\coloneqq a_i\cdot\hat u$ and
  $\alpha^\perp_i\coloneqq \sqrt{|a_i|^2-\alpha_i^2}$ for $1\le i\le n^2$.
  Compute the energy of $b_A$, i.e., $e_A\coloneqq |b_A|^2/n$, which steps~5--6
  will use.

  \emph{Analysis:} The distribution of $\xi$ is invariant under all
  rotations, so the distribution of $\tilde x$ is invariant under the
  rotations that fix $\hat u$.
  \begin{lemma}\label{lem:height}
    Given $\Dd$, assume $\chat\ne0$ and $\tau_0>0$, and let $s$ be the height of
    the proxy input $\tilde x$. Then $\tilde x=s\hat u+\sqrt{1-s^2}\,\chi$, with
    $\chi$ uniform on the unit sphere orthogonal to $\hat u$ and independent of
    $s$. So, given $s$, the first-layer pre-activation $a_i\cdot\tilde x$ has
    the distribution of $s\alpha_i+\sqrt{1-s^2}\,\alpha^\perp_iV$ for each
    $i$, with $V$ a coordinate of a uniform point of
    $\SSS^{n-2}$~\leanmark{1}.
  \end{lemma}
  \begin{proof}
  The lemma rests on the polar decomposition of a standard Gaussian vector
  $g$: $g/|g|$ is uniform on the unit sphere and independent of $|g|$
  \citep[Exercise~3.3.7]{vershynin-2018-hdp}. Given $\Dd$, the vectors
  $\chat$ and $\hat u=\chat/|\chat|$ and the scale $\tau_0$ are fixed, and
  $\xi$ is independent of $\Dd$. Write $\xi=\xi_1\hat u+\xi^\perp$ with
  $\xi_1\coloneqq\xi\cdot\hat u$. Then $\xi_1\sim\Nn(0,1)$, $\xi^\perp$ is a
  standard Gaussian vector in the hyperplane orthogonal to $\hat u$, and
  $\xi_1$ and $\xi^\perp$ are independent. Let $\chi\coloneqq\xi^\perp/|\xi^\perp|$,
  which is defined almost surely since $n\ge2$. By the polar decomposition
  applied to $\xi^\perp$, $\chi$ is uniform on the unit sphere orthogonal to
  $\hat u$ and independent of $(\xi_1,|\xi^\perp|)$. Since
  $\chat+\tau_0\xi=(|\chat|+\tau_0\xi_1)\hat u+\tau_0|\xi^\perp|\,\chi$, the
  height $s$ is a function of $(\xi_1,|\xi^\perp|)$, and so $\chi$ is
  independent of $s$. The coefficient of $\chi$ in $\tilde x$ is nonnegative,
  and $|\tilde x|=1$, so the coefficient is $\sqrt{1-s^2}$. Finally,
  $a_i\cdot \chi=a_i^\perp\cdot \chi$, with $a_i^\perp\coloneqq a_i-\alpha_i\hat u$
  and $|a_i^\perp|=\alpha^\perp_i$. Since the distribution of $\chi$ is invariant
  under the rotations that fix $\hat u$, $a_i^\perp\cdot \chi$ has the
  distribution of $\alpha^\perp_iV$. Hence $a_i\cdot\tilde x$ has the
  distribution of $s\alpha_i+\sqrt{1-s^2}\,\alpha^\perp_iV$ given $s$, as
  claimed.
  \end{proof}
  \begin{leanitems}
  \leanitem[1]{Part161Lat,Part272Reduce}.
  \end{leanitems}
  The height of $\tilde x$ lies near $s_c$, with spread about $\tau_0/\rho$,
  and the nodes span six spreads on each side of $s_c$. $\lambda$ and $S$ vary
  little among proxy inputs of equal height. Step~6 will bound the variation
  of $\lambda$ and $S$.

  \emph{Cost:} $4n^3+O(n^2)$ operations. Almost all of these operations compute the products
  $a_i\cdot\hat u$ and $|a_i|^2$.

\item \textbf{Tables by height.} So that step~6 can estimate $\lambda$ and
  $S$, this step tabulates two quantities of the agent's last hidden vector $h(z)$: the
  alignment of $h(z)$ with $b_A$ and the energy of $h(z)$, i.e.,
  \[f(z)\coloneqq \frac{\langle h(z),b_A\rangle}{\sqrt n\,|b_A|},\qquad
  \omega(z)\coloneqq \frac{|h(z)|^2}{n}\qquad\text{for }z\in\RR^n.\]
  This step estimates the averages of $f$ and $\omega$ over proxy inputs whose
  height is a node.

  \emph{Computation:} Let $J_z\coloneqq \lceil n^{1.5}\rceil-1$, and let $V$ be
  as in step~4. For each node $g$ and $1\le i\le n^2$, let $X_{g,i}$ be
  first-layer unit $i$ at height $s_g$, i.e.,
  $X_{g,i}\coloneqq \relu(s_g\alpha_i+\sqrt{1-s_g^2}\,\alpha^\perp_iV)$.
  Compute estimates $\bar h_{g,i}$ and $\bar h^{(2)}_{g,i}$ of $\EE_VX_{g,i}$
  and $\EE_VX_{g,i}^2$. For even $n$, use the exact closed forms of these expectations. For odd
  $n$, use a finite arithmetic rule whose error the Numerical certificates
  will bound~\leanmark{1}. Estimate the spread of the second-layer
  pre-activations, i.e.,
  $\varsigma_g\coloneqq \sqrt{\nu_A(\sum_i\bar h^{(2)}_{g,i}-|\bar h_g|^2)}$.
  For $1\le j\le n^2$, let
  $\phi_{g,j}(z)\coloneqq \relu((A_2\bar h_g)_j+\varsigma_gz)$, and estimate
  the first two moments of last-hidden-layer unit $j$ by
  $\beta_{g,j}\coloneqq B_{J_z}[\phi_{g,j}]$ and
  $\beta^{(2)}_{g,j}\coloneqq B_{J_z}[\phi_{g,j}^2]$. Store the estimates
  $\tilde f_g$ and $\tilde\omega_g$ of the averages of $f$ and $\omega$ at
  height $s_g$, i.e.,
  \[\tilde f_g\coloneqq \frac{1}{\sqrt n\,|b_A|}\sum_j\beta_{g,j}b_{A,j},\qquad
  \tilde\omega_g\coloneqq \frac1n\sum_j\beta^{(2)}_{g,j}.\]
  \begin{leanitems}
  \leanitem[1]{Part272Stated,Part272Even,Part278Err}.
  \end{leanitems}

  \emph{Analysis:} This step tabulates $f$ and $\omega$ because the
  coefficient and the noise variance are functions of them.
  \begin{lemma}\label{lem:fomega}
    The coefficient and the noise variance of Lemma~\ref{lem:agentout} are
    functions of $f$, $\omega$ and $e_A$, i.e.,
    $\lambda(z)=f(z)/\sqrt{e_A}$ and $S(z)=\gamma_A(\omega(z)-f(z)^2)$ for
    every $z\in\RR^n$, where $\gamma_A\coloneqq n\nu_A=2/n$~\leanmark{1}.
  \end{lemma}
  \begin{proof}
  Fix $z\in\RR^n$, and write $h$, $h^\perp$, $\lambda$, $S$, $f$ and $\omega$
  for their values at $z$. First let $b_A\ne0$. Then
  $\sqrt{e_A}=|b_A|/\sqrt n$, so
  $f/\sqrt{e_A}=\langle h,b_A\rangle/|b_A|^2=\lambda$. Since
  $h=\lambda b_A+h^\perp$ with $h^\perp$ orthogonal to $b_A$,
  $|h^\perp|^2=|h|^2-\lambda^2|b_A|^2
  =|h|^2-\langle h,b_A\rangle^2/|b_A|^2=n(\omega-f^2)$. So
  $S=\nu_A|h^\perp|^2=n\nu_A(\omega-f^2)=\gamma_A(\omega-f^2)$. Now let
  $b_A=0$. Then $f/\sqrt{e_A}=0=\lambda$, since $v/0=0$; so $h^\perp=h$, and
  $S=\nu_A|h|^2=\gamma_A\omega=\gamma_A(\omega-f^2)$. In both cases,
  $\nu_A=2/n^2$, so $\gamma_A=2/n$. Hence the two identities hold for every
  $z$, as claimed.
  \end{proof}
  \begin{leanitems}
  \leanitem[1]{Part169Var,Part169Red}. Lean proves the identity for
  $\lambda$ in the form $f(z)=\sqrt{e_A}\,\lambda(z)$ as \texttt{fX\_eq\_lam},
  and the identity for $S$ as \texttt{sA\_eq}, which assumes $b_A\ne0$.
  \end{leanitems}
  By
  Lemma~\ref{lem:height}, when $n$ is even, $\bar h_{g,i}$ and
  $\bar h^{(2)}_{g,i}$ are the first two moments of first-layer unit $i$ over
  proxy inputs of height $s_g$; when $n$ is odd, $\bar h_{g,i}$ and $\bar h^{(2)}_{g,i}$ estimate these moments.
  The second layer is estimated as in step~3.
  So $\tilde f_g$ and $\tilde\omega_g$ are designed to approximate the averages
  of $f$ and $\omega$ over proxy inputs of height $s_g$.

  \emph{Cost:} $1024n^5+O(n^{4.5})$ operations, i.e., $512n$ nodes at
  $2n^4+O(n^{3.5})$ each. Almost all of these operations compute the products $A_2\bar h_g$. At
  $n=8{,}000$, the leading term $1024n^5$ is $8\cdot10^{-6}$ times $2n^7$.
  The cost model in Lean does not charge for the weights of $B_{J_z}$ or for
  the coefficients of the arithmetic rules for $\bar h_{g,i}$ and
  $\bar h^{(2)}_{g,i}$, such as $(-1)^q\binom Pq/(2q+j+1)$ with
  $P\coloneqq (n-4)/2$ for even $n$. Computing these weights and coefficients once would take about
  $3J_z+O(n)$ operations, i.e., about $2\cdot10^6$ at $n=8{,}000$, less than
  $10^{-16}$ times $1024n^5$.

\item \textbf{Proxy outputs.} This step draws the heights of $K$ proxy
  inputs. For each proxy input, this step draws a proxy output from the
  distribution of Lemma~\ref{lem:agentout}, with $\lambda$ and $S$ replaced by their tabulated
  estimates.

  \emph{Computation:} To keep the noise variance positive, $\Est$ uses a
  floor $S_0$ for $\tilde\omega_g-\tilde f_g^2$: let $S_0\coloneqq 0.17$ if $n<11{,}883$ and $S_0\coloneqq 0.2$
  otherwise. Read the seed as $r=(r_k)_{k=1}^K$ with $K\coloneqq 10^{12}$ and
  independent standard Gaussian vectors
  $r_k=(Z_k,G_{k,1},\dots,G_{k,n-1},\eta_k)\in\RR\times\RR^{n-1}\times\RR^n$.
  For each $k=1,\dots,K$, compute the height $s_k$ of the $k$-th proxy
  input from the coordinate $w_k$ of the unnormalized proxy input along
  $\hat u$, i.e.,
  \[w_k\coloneqq |\chat|+\tau_0Z_k,\qquad
  s_k\coloneqq \frac{w_k}{\sqrt{w_k^2+\tau_0^2\sum_iG_{k,i}^2}}.\]
  Let $g_k\in\{0,\dots,512n-1\}$ index the node nearest to $s_k$, the larger
  at a tie. Compute the coefficient $\tilde\lambda_k$ and the noise variance $\tilde S_k$
  from $\tilde f_{g_k}$ and $\tilde\omega_{g_k}$ by Lemma~\ref{lem:fomega}, and form the proxy output $y_k$, i.e.,
  \[\tilde\lambda_k\coloneqq \frac{\tilde f_{g_k}}{\sqrt{e_A}},\qquad
  \tilde S_k\coloneqq \gamma_A\max(S_0,\tilde\omega_{g_k}-\tilde f_{g_k}^2),\qquad
  y_k\coloneqq \tilde\lambda_kc_A+\sqrt{\tilde S_k}\,\eta_k.\]

  \emph{Analysis:} The vector $(w_k,\tau_0G_{k,1},\dots,\tau_0G_{k,n-1})$ has
  the distribution of the unnormalized proxy input $\chat+\tau_0\xi$ in
  coordinates whose first axis is $\hat u$.
  \begin{lemma}\label{lem:seedheight}
    Given $\Dd$, if $\chat\ne0$ or $\tau_0>0$, then $s_k$ has exactly the
    distribution of the height $\tilde x\cdot\hat u$~\leanmark{1}.
  \end{lemma}
  \begin{proof}
  The lemma rests on the rotation invariance of the standard Gaussian
  distribution \citep[Proposition~3.3.2]{vershynin-2018-hdp}. Given $\Dd$,
  the vectors $\chat$ and $\hat u$ and the scale $\tau_0$ are fixed, and
  $\chat=|\chat|\hat u$, also when $\chat=0$. Let $\mathcal Q$ be an orthogonal
  $n\times n$ matrix with first column $\hat u$, and let
  $\xi'\coloneqq \mathcal Q^\top\xi$. Then $\xi'\sim\Nn(0,I)$, and
  $\mathcal Q^\top(\chat+\tau_0\xi)=(|\chat|+\tau_0\xi'_1,\tau_0\xi'_2,\dots,\tau_0\xi'_n)$.
  Since $\mathcal Q^\top$ preserves lengths and maps $\hat u$ to $e_1$,
  \[\tilde x\cdot\hat u=\frac{|\chat|+\tau_0\xi'_1}
  {\sqrt{(|\chat|+\tau_0\xi'_1)^2+\tau_0^2\sum_{i=2}^n\xi_i'^2}}.\]
  The right side is the formula for $s_k$, with $\xi'$ in place
  of $(Z_k,G_{k,1},\dots,G_{k,n-1})$. Both of these vectors have the
  distribution $\Nn(0,I)$. If $\tau_0>0$, the denominator is $0$ only on an
  event of probability $0$; if $\tau_0=0$, then $\chat\ne0$ and both $\tilde x\cdot\hat u$ and $s_k$
  equal $1$. Hence $s_k$ has the distribution of $\tilde x\cdot\hat u$, as
  claimed.
  \end{proof}
  \begin{leanitems}
  \leanitem[1]{Part161Lat}. Lean proves the lemma as part of
  \texttt{nuStand\_latitude}, which assumes $\chat\ne0$ and $\tau_0>0$.
  \end{leanitems}
  Flooring $\tilde\omega_{g_k}-\tilde f_{g_k}^2$ at $S_0$ does not increase
  the absolute error of this difference as an estimate of a target value that
  is at least $S_0$.
  \begin{lemma}\label{lem:floor}
    Flooring never increases the distance to a target value $Y\ge S_0$,
    i.e., $|\max(S_0,v)-Y|\le|v-Y|$ for reals $v$ and $Y\ge S_0$~\leanmark{1}.
  \end{lemma}
  \begin{proof}
  If $v\ge S_0$, then $\max(S_0,v)=v$, and $|\max(S_0,v)-Y|=|v-Y|$. If
  $v<S_0$, then $\max(S_0,v)=S_0$ and $v<S_0\le Y$, so
  $|\max(S_0,v)-Y|=Y-S_0<Y-v=|v-Y|$. Hence $|\max(S_0,v)-Y|\le|v-Y|$ for
  every real $v$, as claimed.
  \end{proof}
  Let $\tilde q$ be the probability that one proxy output fails, i.e.,
  $\tilde q(\Dd)\coloneqq \PP_r[F(y_1)>t\mid\Dd]$. For probability
  densities $p_1$ and $p_2$, let the \emph{Hellinger distance} be
  $H(p_1,p_2)\coloneqq\bigl(\int(\sqrt{p_1}-\sqrt{p_2})^2\bigr)^{1/2}$, with
  no factor $1/2$ inside the root. Two Gaussians
  $\Nn(\lambda c_A,SI)$ with nearby $(\lambda,S)$ are close in Hellinger
  distance. Among proxy inputs of equal height, $\lambda$ and $\sqrt S$ deviate
  little from their conditional means; Lean bounds moments of these
  deviations, averaged over the height, and the mixture comparison also
  accounts for exceptional directions~\leanmark{2}. A Hellinger bound on
  the distribution of the output bounds the difference of the square roots of
  the failure probabilities. This gives the following lemma.
  \begin{lemma}\label{lem:da}
    For every $t\in\RR$ and on the good set $\Gg$ of ``Bounds over the residual
    parts'', the average of $\tilde p$ over $\Omega_A$ is close to
    $\tilde q(\Dd)$ on the square-root scale, i.e.,
    \[\bigl|\sqrt{\EE_{\Omega_A}\tilde p}-\sqrt{\tilde q(\Dd)}\bigr|\le\Delta_A~\leanmark{2},\]
    where $\Delta_A>0$ is an explicit real depending only on $n$, given in the
    Numerical certificates.
  \end{lemma}
  \begin{proof}[Proof sketch]
  The full proof is in Lean. Fix $t$ and $\Dd\in\Gg$, and let
  $B\coloneqq\{y\in\RR^n:F(y)>t\}$. The proof combines three standard
  inequalities. The first inequality is that
  \begin{equation}\label{eq:hellevent}
  \bigl(\sqrt{P_1(B)}-\sqrt{P_2(B)}\bigr)^2\le H(p_1,p_2)^2
  \end{equation}
  for distributions $P_1$ and $P_2$ with densities $p_1$ and $p_2$, by the
  Cauchy--Schwarz inequality on $B$ and on $\RR^n\setminus B$. The second
  inequality is a bound on the Hellinger distance between two isotropic
  Gaussians. By the closed form of \citet[Section~2]{dmr-2022}, multiplied
  by $2$ because $H$ has no factor $1/2$ inside the root, with $\varrho\coloneqq S_1/S_2$,
  \begin{equation}\label{eq:hellgauss}
  \begin{aligned}
  H\bigl(\Nn(m_1,S_1I),\Nn(m_2,S_2I)\bigr)^2
  &=2-2\Bigl(\frac{2\sqrt{S_1S_2}}{S_1+S_2}\Bigr)^{n/2}
  e^{-|m_1-m_2|^2/(4(S_1+S_2))}\\
  &\le\frac{n(\varrho-1)^2}{8\varrho}+\frac{|m_1-m_2|^2}{2(S_1+S_2)}
  \end{aligned}
  \end{equation}
  for $S_1,S_2>0$, where the inequality follows from $1-y\le-\ln y$ and
  $\ln y\le y-1$ for $y>0$. The third inequality is that
  \begin{equation}\label{eq:hellmix}
  \bigl(\sqrt{\EE X}-\sqrt{\EE Y}\bigr)^2\le\EE\bigl(\sqrt X-\sqrt Y\bigr)^2
  \end{equation}
  for random variables $X,Y$ with values in $[0,1]$, by the Cauchy--Schwarz
  inequality.

  By Lemma~\ref{lem:agentout}, $\EE_{\Omega_A}\tilde p$ is the average over
  $\tilde x$ of $\Nn(\lambda(\tilde x)c_A,S(\tilde x)I)(B)$. By
  Lemma~\ref{lem:seedheight}, $\tilde q(\Dd)$ is the average over the height
  $s$ of $\tilde x$ of the same Gaussian probability, with the tabulated
  coefficient and noise variance at the node nearest $s$ in place of
  $\lambda(\tilde x)$ and $S(\tilde x)$. The Lean proof bounds the difference
  of the square roots of these two averages by a sum of four terms. The first
  term is the fluctuation of $(\lambda,S)$ among proxy inputs of equal height,
  bounded through the $\chi^2$ divergence of a Gaussian mixture. The second
  term is the bias of the tables against the averages of $(\lambda,S)$ at
  each height, bounded by \eqref{eq:hellevent}--\eqref{eq:hellmix} and by
  Lemma~\ref{lem:floor}. The third term comes from the heights outside the
  range of the nodes. The fourth term is the error of the rule of step~5 for
  odd $n$, at most $10^{-20}$. The Numerical certificates will bound the sum of
  the four terms by $\Delta_A$, as claimed.
  \end{proof}
  \begin{leanitems}
  \leanitem[1]{Part289Floor}.
  \leanitem[2]{Part157Repl,Part161Root,Part161Gauss,Part161Red,Part161Fl,Part281Odd,Part289Root,Part293Root,Part293Arithmetic};
  \href{\leanrepo}{\leanstepE}; \href{\leanrepo}{\leanstepJ}.
  \end{leanitems}

  \emph{Cost:} at most $K(5n+10\lceil\log_2(512n)\rceil+27)$ operations, i.e., about
  $4\cdot10^{16}$ at $n=8{,}000$. The cost model in Lean charges nothing
  for the seed. Generating the seed's $2nK$ Gaussians from random bits would
  take about $200Kn$ operations at an estimated $100$ operations per Gaussian
  (e.g., by the Box--Muller method with series for $\ln$, $\cos$ and $\sin$),
  i.e., about $1.6\cdot10^{18}$ at $n=8{,}000$. This
  generation cost is about $40$ times the charged cost
  of this step, but less than $4\cdot10^{-10}$ times $2n^7$, and less than
  $2\cdot10^{-6}$ times the cost of step~7.

\item \textbf{The failure fraction and the output.} This step evaluates $F$ on
  the proxy outputs. This step then scales and clips the fraction of proxy
  outputs that fail, which gives the output $U$.

  \emph{Computation:} Let $R$ be the fraction of proxy outputs that fail,
  i.e.,
  \[R\coloneqq K^{-1}\sum_{k=1}^K\mathbf 1[F(y_k)>t].\]
  Let $\Gamma\coloneqq 2.2524763228$. Output the scaled and clipped
  fraction, i.e.,
  \[U\coloneqq \min(1,\max(1/100,\Gamma R)).\]

  \emph{Analysis:} $\Est$ succeeds whenever $R$ lies in the \emph{window}
  $[p/\Gamma,\ \max(10p,1/100)/\Gamma]$. The lower end matters only when
  $p>1/100$, since otherwise $U\ge1/100\ge p$ by the definition of $U$. Since $1/\sqrt\Gamma<0.6664$
  and $\sqrt{10/\Gamma}>2.1070$, the window gives the following lemma on the
  square-root scale.
  \begin{lemma}\label{lem:window}
    If $p<1/1000$, $\Est$ succeeds whenever $R\le1/(100\Gamma)$. If
    $p\ge1/1000$, $\Est$ succeeds whenever
    $0.667\sqrt p\le\sqrt R\le2.107\sqrt p$~\leanmark{1}.
  \end{lemma}
  \begin{proof}
  Recall that $0\le p\le1$ and $R\ge0$. First let $p<1/1000$ and
  $R\le1/(100\Gamma)$. Then $\Gamma R\le1/100$, so $U=1/100$, and
  $p<U\le\max(10p,1/100)$. Now let $p\ge1/1000$ and
  $0.667\sqrt p\le\sqrt R\le2.107\sqrt p$. Since $0.667^2\Gamma>1.0021$ and
  $2.107^2\Gamma<9.9998$, squaring and multiplying by $\Gamma$ give $p\le\Gamma R\le10p$. So
  $U\ge\min(1,\Gamma R)\ge\min(1,p)=p$, and
  $U\le\max(1/100,\Gamma R)\le\max(1/100,10p)$. Hence $\Est$ succeeds in both
  cases, as claimed.
  \end{proof}
  \begin{leanitems}
  \leanitem[1]{Part156Out}. Lean proves, as \texttt{not\_success}, the
  contrapositive of a slightly stronger form of the statement that $\Est$ succeeds whenever $R$ lies in
  the window, for a factor $a>0$ in place
  of $\Gamma$, and with the window stated for a value $Q$ near $R$ on the
  square-root scale, i.e., $|\sqrt R-\sqrt Q|$ at most a given error. The
  lemma follows with $a=\Gamma$, $Q=R$ and error $0$, by the inequalities for $0.667^2\Gamma$ and
  $2.107^2\Gamma$ in its proof.
  \end{leanitems}
  Given $\Dd$, the $K$ proxy outputs fail independently, each with
  probability $\tilde q(\Dd)$. A Chernoff bound then keeps $\sqrt R$ near
  $\sqrt{\tilde q}$.
  \begin{lemma}\label{lem:chernoff}
    If, given $\Dd$, $KR\sim\mathrm{Bin}(K,\tilde q(\Dd))$, then, for
    every real $\Delta_s>0$, $\sqrt R$ is close to $\sqrt{\tilde q(\Dd)}$,
    i.e.,
    $\PP_r(|\sqrt R-\sqrt{\tilde q(\Dd)}|>\Delta_s\mid\Dd)\le2e^{-K\Delta_s^2}$~\leanmark{1}.
  \end{lemma}
  \begin{proof}
  Lemma~\ref{lem:chernoff} is standard; see \citet[Theorems~3--4]{okamoto-1958},
  which give $e^{-2K\Delta_s^2}+e^{-K\Delta_s^2}$ in place of
  $2e^{-K\Delta_s^2}$. Condition on $\Dd$, and let $q\coloneqq\tilde q(\Dd)$.
  If $q\in\{0,1\}$, then $R=q$ almost surely, and the probability is $0$. Let
  $0<q<1$, and let $D(z\Vert q)\coloneqq z\ln(z/q)+(1-z)\ln((1-z)/(1-q))$ for
  $0<z<1$. By the Chernoff bound \citep[Lemma~1]{okamoto-1958},
  $\PP_r(R\ge z\mid\Dd)\le e^{-KD(z\Vert q)}$ for $z\ge q$, and
  $\PP_r(R\le z\mid\Dd)\le e^{-KD(z\Vert q)}$ for $z\le q$. By
  $\ln y\le y-1$ at $y=\sqrt{q/z}$ and at $y=\sqrt{(1-q)/(1-z)}$,
  \[-D(z\Vert q)\le2\bigl(\sqrt{zq}-z\bigr)+2\bigl(\sqrt{(1-z)(1-q)}-(1-z)\bigr)
  =-(\sqrt z-\sqrt q)^2-(\sqrt{1-z}-\sqrt{1-q})^2,\]
  so $D(z\Vert q)\ge(\sqrt z-\sqrt q)^2$. The event
  $\sqrt R>\sqrt q+\Delta_s$ is empty if $(\sqrt q+\Delta_s)^2\ge1$, and
  otherwise lies in $\{R\ge z_+\}$ with $z_+\coloneqq(\sqrt q+\Delta_s)^2\in(q,1)$.
  The event $\sqrt R<\sqrt q-\Delta_s$ is empty if $\sqrt q\le\Delta_s$, and
  otherwise lies in $\{R\le z_-\}$ with $z_-\coloneqq(\sqrt q-\Delta_s)^2\in(0,q)$.
  For $z=z_+$ and for $z=z_-$, $(\sqrt z-\sqrt q)^2=\Delta_s^2$, so each event has
  probability at most $e^{-K\Delta_s^2}$. Hence the union of the two events
  has probability at most $2e^{-K\Delta_s^2}$, as claimed.
  \end{proof}
  \begin{leanitems}
  \leanitem[1]{Part156Samp,Part232World}.
  \end{leanitems}
  With $K=10^{12}$, the bound of Lemma~\ref{lem:chernoff} is $2e^{-100}$ at
  $\Delta_s=10^{-5}$, the value Lean uses for $n\ge11{,}883$, and
  $2e^{-25}<10^{-9}$ at $\Delta_s=5\cdot10^{-6}$, the value Lean uses for
  $8{,}000\le n\le11{,}882$. Errors on the square-root scale
  suit the window, as follows. For $0\le\Delta\le\sqrt p$, if
  $|\sqrt R-\sqrt p|\le\Delta$, then
  $(1-\Delta/\sqrt p)^2\le R/p\le(1+\Delta/\sqrt p)^2$. Also,
  $\sqrt p\ge0.0316$ wherever $\max(10p,1/100)=10p$, i.e.,
  $p\ge1/1000$. Hence, if $p\ge1/1000$ and $|\sqrt R-\sqrt p|\le0.01$, then
  $0.467\le R/p\le1.733$, inside $[1/\Gamma,10/\Gamma]$, the window's range for
  $R/p$, as claimed.

  \emph{Cost:} $K(2m_F^2+2nm_F+O(m_F))$ operations, i.e., one evaluation of $F$
  on each of the $K$ proxy outputs, where $m_F\coloneqq \lceil n^{1.5}\rceil$ is the hidden
  width of $F$.
\end{enumerate}

\subsection{Completion of the proof}

\paragraph{Bounds over the residual parts.}
This paragraph proves three bounds: the $\Delta_W$ link of the chain~\eqref{eq:chain}, and
the concentration bounds for the two links marked ``conc.''. Given $\Dd$, the
matrices $\Omega_W$ and $\Omega_A$ and the seed $r$ are independent and
Gaussian, so $p$ is a function of
$(\Dd,\Omega_W,\Omega_A)$ and $\tilde p$ is a function of $(\Dd,\Omega_A)$.
All three bounds hold on a good set $\Gg$ of values of $\Dd$. To define $\Gg$,
let $\delta_{\mathrm{hid}},\mu_W,\delta_{\mathrm{ext}},\Delta_W>0$ be explicit
reals that depend only on $n$; for an example, the Numerical certificates will
give $\Delta_W$ and a bound on $\delta_{\mathrm{hid}}+\mu_W/\Delta_W^2+\delta_{\mathrm{ext}}$. For each $u$, let $H_u$ be the Hellinger distance of step~6 between
the noisy copy of $\chat$ in Lemma~\ref{lem:noisy} and the noisy copy of
$\chat$ that defines the proxy input, i.e.,
$H_u\coloneqq H\bigl(\Nn(\chat,\tau_W(u)^2I),\Nn(\chat,\tau_0^2I)\bigr)$ if
$\lambda_W(u)>0$, $\tau_W(u)>0$ and $\tau_0>0$, and $H_u\coloneqq\sqrt2$
otherwise. Let $\bar\Delta$ be the root mean square of $H_u$, i.e.,
$\bar\Delta(\Dd)\coloneqq(\EE_uH_u^2)^{1/2}$, a function of $\Dd$. Let $\Hh$
be a set of values of $\Dd$ and $\Ee$ a further set of values of $\Dd$, such
that
\[\PP(\Dd\notin\Hh)\le\delta_{\mathrm{hid}},\qquad
\EE[\mathbf 1_{\Hh}\bar\Delta^2]\le\mu_W,\qquad
\PP(\Dd\in\Hh\setminus\Ee)\le\delta_{\mathrm{ext}}.\]
On $\Hh$, the world's hidden layers satisfy explicit bounds, which only the Lean proof states.
Let the good set be $\Gg\coloneqq \Hh\cap\Ee\cap\{\bar\Delta\le\Delta_W\}$.

The concentration bounds depend on the directions of the vectors
$h^\perp(\tilde x)$. Let $\zeta$ be the direction of $h^\perp$, i.e.,
$\zeta\coloneqq h^\perp/|h^\perp|$ if $h^\perp\ne0$ and $\zeta\coloneqq0$ otherwise, and let the frame constant of $\zeta$ be
$L_A\coloneqq \|\EE_{\tilde x}\zeta(\tilde x)\zeta(\tilde x)^\top\|_{\mathrm{op}}$.
Let $L_W$ be the same for the world, i.e., with $h_W^\perp$ and $u$ in place
of $h^\perp$ and $\tilde x$. For real $L>0$ and $\sigma\in\RR$, let
$\chi_L(\sigma)\coloneqq (e^{\sigma L}-1)/L$. Lemma~\ref{lem:conc} below will bound the moment
generating functions of $p$ and $\tilde p$, by Lemma~\ref{lem:gbl}, a Gaussian
form of the geometric Brascamp--Lieb inequality
\citep[Proposition~2.8]{bcct-2008}. Lemma~\ref{lem:frame}
will bound $L_A$, by the Gaussian Poincar\'e inequality and a Schur test on
pairs of proxy inputs.
\begin{lemma}\label{lem:dw}
  On $\Gg$, $\tilde p$ is close to the average of $p$ over $\Omega_W$, i.e.,
  \[\bigl|\sqrt{\EE_{\Omega_W}p}-\sqrt{\tilde p}\bigr|\le\Delta_W\ \text{for every }\Omega_A,
  \quad\text{and}\quad\PP(\Dd\notin\Gg)\le\delta_{\mathrm{hid}}+\mu_W/\Delta_W^2+\delta_{\mathrm{ext}}.~\leanmark{1}\]
\end{lemma}
\begin{proof}
The first claim needs only $\bar\Delta\le\Delta_W$. Fix $t$, $\Omega_A$, and $\Dd$ with
$\bar\Delta(\Dd)\le\Delta_W$, and let
$T\coloneqq\{z\in\SSS^{n-1}:F(A(z))>t\}$. For each $u$, let $P_u$ be the
probability over $\Omega_W$ that $x\in T$. By Fubini's theorem,
$\EE_{\Omega_W}p=\EE_uP_u$. Fix $u$, and let $f$ and $g$ be the densities of
$\Nn(\chat,\tau_W(u)^2I)$ and $\Nn(\chat,\tau_0^2I)$, respectively. If $H_u=\sqrt2$, then
$(\sqrt{P_u}-\sqrt{\tilde p})^2\le1\le H_u^2$, since
$P_u,\tilde p\in[0,1]$. Otherwise, by Lemma~\ref{lem:noisy} and the
definition of $\tilde x$, $P_u=\int_Sf$ and $\tilde p=\int_Sg$ for the set
$S$ of vectors whose direction lies in $T$, up to a null set. By the
Cauchy--Schwarz inequality, $\int_S\sqrt{fg}\le\sqrt{P_u\tilde p}$, so
$(\sqrt{P_u}-\sqrt{\tilde p})^2\le\int_S(\sqrt f-\sqrt g)^2\le H_u^2$. By the
Cauchy--Schwarz inequality again,
$\bigl(\sqrt{\EE_uP_u}-\sqrt{\tilde p}\bigr)^2\le\EE_u\bigl(\sqrt{P_u}-\sqrt{\tilde p}\bigr)^2
\le\EE_uH_u^2=\bar\Delta^2$. Hence
$|\sqrt{\EE_{\Omega_W}p}-\sqrt{\tilde p}|\le\bar\Delta\le\Delta_W$, and in
particular on $\Gg$. For the second claim, $\Dd\notin\Gg$ only if
$\Dd\notin\Hh$, or $\Dd\in\Hh\setminus\Ee$, or $\Dd\in\Hh$ and
$\bar\Delta>\Delta_W$. By Markov's inequality, the last event has probability
at most $\EE[\mathbf 1_\Hh\bar\Delta^2]/\Delta_W^2\le\mu_W/\Delta_W^2$. Hence,
by a union bound,
$\PP(\Dd\notin\Gg)\le\delta_{\mathrm{hid}}+\mu_W/\Delta_W^2+\delta_{\mathrm{ext}}$,
as claimed.
\end{proof}
\begin{lemma}[Gaussian geometric Brascamp--Lieb inequality]\label{lem:gbl}
  Let $\mu$ be a Borel probability measure on a Euclidean space
  $\mathcal Z$, let
  $\zeta:\mathcal Z\to\RR^m$ be measurable with $|\zeta(z)|\in\{0,1\}$ for every $z$,
  and let $L>0$ with $\int\zeta\zeta^\top d\mu\preceq L\,I$. Let $G$ be an
  $n\times m$ matrix of independent $\Nn(0,1)$ entries, and let
  $Z\sim\Nn(0,I)$ in $\RR^n$. Let $0<a\le b$, and let $(f_z)_{z\in\mathcal Z}$ be
  a family of functions $f_z:\RR^n\to[a,b]$, jointly measurable in $z$ and in the argument of $f_z$, such
  that $f_z$ is constant for every $z$ with $\zeta(z)=0$. Then
  \[\EE_G\exp\Bigl(L^{-1}\!\int\ln f_z(G\zeta(z))\,d\mu(z)\Bigr)
  \le\exp\Bigl(L^{-1}\!\int\ln\EE_Zf_z(Z)\,d\mu(z)\Bigr)~\leanmark{2}.\]
\end{lemma}
\begin{proof}
Lemma~\ref{lem:gbl} follows from the Bou\'e--Dupuis formula
\citep[Proposition~1.1]{lehec-2014}, by the argument with which
\citet[Section~2]{lehec-2014} proves the Brascamp--Lieb inequality.
\end{proof}
\begin{lemma}\label{lem:conc}
  Given $\Dd$, let $L'_W\ge L_W$ and $L'_A\ge L_A$ be positive reals. Then
  $p$ and $\tilde p$ concentrate around their averages over $\Omega_W$ and
  $\Omega_A$, respectively, i.e.,
  \[\EE_{\Omega_W}e^{\sigma p}\le e^{\chi_{L'_W}(\sigma)\,\EE_{\Omega_W}p},\qquad
  \EE_{\Omega_A}e^{\sigma\tilde p}\le e^{\chi_{L'_A}(\sigma)\,\EE_{\Omega_A}\tilde p}\]
  for every $\sigma\in\RR$ and, in the first bound, every
  $\Omega_A$~\leanmark{2}.
\end{lemma}
\begin{proof}
Consider $\tilde p$. The bound for $p$ is the same, for each fixed
$\Omega_A$, with $u$, $h_W^\perp$, $\nu_W$, $L'_W$ and $W_3$ in place of
$\tilde x$, $h^\perp$, $\nu_A$, $L'_A$ and $A_3$. Let $L\coloneqq L'_A$, and
let $\mu$ be the distribution of $\tilde x$, so that
$\int\zeta\zeta^\top d\mu\preceq L\,I$. As in the proof of
Lemma~\ref{lem:agentout}, given $\Dd$, the matrix $\Omega_A$ has the
distribution of $\sqrt{\nu_A}\,G\Pi_{b_A}$, with $G$ an $n\times n^2$ matrix of
independent $\Nn(0,1)$ entries. So
$\Omega_Ah^\perp(z)=\sqrt{\nu_A}\,|h^\perp(z)|\,G\zeta(z)$ for every $z$,
jointly in distribution, and $\tilde p=\int\phi_z(G\zeta(z))\,d\mu(z)$, with
$\phi_z(y)\coloneqq\mathbf 1[F(\lambda(z)c_A+\sqrt{\nu_A}\,|h^\perp(z)|\,y)>t]$.
The function $\phi_z$ is constant where $\zeta(z)=0$. Apply
Lemma~\ref{lem:gbl} with $m=n^2$ and $f_z\coloneqq e^{\sigma L\phi_z}$. The
left side of Lemma~\ref{lem:gbl} is
$\EE_{\Omega_A}e^{\sigma\tilde p}$. Since $\phi_z\in\{0,1\}$,
$\EE_Zf_z(Z)=1+(e^{\sigma L}-1)\EE_Z\phi_z(Z)$, and $\ln(1+y)\le y$ for
$y>-1$. So the right side of Lemma~\ref{lem:gbl} is at most
$\exp\bigl(\chi_L(\sigma)\int\EE_Z\phi_z(Z)\,d\mu(z)\bigr)$. Where
$\zeta(z)\ne0$, $G\zeta(z)\sim\Nn(0,I)$; so, by Fubini's theorem,
$\int\EE_Z\phi_z(Z)\,d\mu(z)=\EE_{\Omega_A}\tilde p$. Hence
$\EE_{\Omega_A}e^{\sigma\tilde p}\le e^{\chi_{L'_A}(\sigma)\,\EE_{\Omega_A}\tilde p}$,
as claimed.
\end{proof}
\begin{lemma}\label{lem:frame}
  On $\Gg$, the frame constants are of order $1/n$, i.e., $L_W\le \vartheta_W/n$
  and $L_A\le \vartheta_A/n$ for explicit reals $\vartheta_W,\vartheta_A>0$
  that depend only on $n$; the Numerical certificates give examples of
  their values~\leanmark{3}.
\end{lemma}
\begin{proof}[Proof sketch for $L_A$]
The full proof is in Lean. Fix $\Dd\in\Gg$, with $\tau_0>0$. Let
$y\coloneqq\chat+\tau_0\xi$, so that $\tilde x=y/|y|$, let
$w\coloneqq\tau_0^2/|y|^2$, let $J(z)$ be the Jacobian of $h$ at $z$, and let
$\Pi_z\coloneqq I-zz^\top$. The sketch has three parts.

The first part of the sketch bounds $L_A$ by a cutoff on $S$. Since $\zeta\zeta^\top=\nu_Ah^\perp h^{\perp\top}/S$
and $\zeta\zeta^\top\preceq I$, for every real $d_0>0$,
\begin{equation}\label{eq:framecut}
L_A\le\frac{\nu_A}{\gamma_Ad_0}\bigl(\|\Pi_{b_A}\operatorname{Cov}_{\tilde x}(h)\,\Pi_{b_A}\|_{\mathrm{op}}
+|\Pi_{b_A}\EE_{\tilde x}h|^2\bigr)+\PP_{\tilde x}(S<\gamma_Ad_0),
\end{equation}
with $\Pi_{b_A}$ as in Lemma~\ref{lem:gauss}.

The second part bounds the covariance by the Gaussian Poincar\'e inequality
\citep[Theorem~3.20]{blm-2013}, which states that
$\operatorname{Var}f(\xi)\le\EE|\nabla f(\xi)|^2$ for continuously
differentiable $f$. Lean proves the inequality for Lipschitz $f$, since $h$ is
Lipschitz but not continuously differentiable. As $y/|y|$ is not Lipschitz
near $y=0$, Lean applies it with $y/\max\{|y|,\upsilon\}$ in place of $y/|y|$,
for $\upsilon>0$, and lets $\upsilon\to0$, using $\EE|y|^{-4}<\infty$. For each $v\in\RR^{n^2}$, the chain rule gives
$\nabla_\xi\langle h(\tilde x),v\rangle=(\tau_0/|y|)\,\Pi_{\tilde x}J(\tilde x)^\top v$
almost everywhere. Hence
\begin{equation}\label{eq:framepoinc}
\nu_A\operatorname{Cov}_{\tilde x}(h)\preceq D,\qquad D\coloneqq\EE_\xi BB^\top,\qquad
B\coloneqq\sqrt{\nu_Aw}\,J(\tilde x)\Pi_{\tilde x}.
\end{equation}

The third part bounds $D$ by pairs of proxy inputs. Let $T$ map $f$ in
$L^2(\xi;\RR^n)$ to $\EE_\xi Bf(\xi)$. Then
$\|D\|_{\mathrm{op}}=\|TT^*\|_{\mathrm{op}}=\sup_{\|f\|_2=1}|Tf|^2$, and $|Tf|^2$
is the average of $f(\xi)^\top B^\top B'f(\xi')$ over two independent proxy
inputs, with $B'$ the value of $B$ at the second input. The Lean proof
bounds this average by the Schur test \citep[Theorem~5.2]{halmos-sunder-1978},
with explicit charges for exceptional pairs of inputs. With this bound
on $\|D\|_{\mathrm{op}}$, \eqref{eq:framepoinc} bounds the covariance term
of~\eqref{eq:framecut}. The Lean proof bounds
the other two terms of~\eqref{eq:framecut} on $\Gg$, and then give
$L_A\le\vartheta_A/n$, as claimed. On the range $8{,}000\le n\le11{,}882$,
they use a variant of~\eqref{eq:framecut}.
\end{proof}
\begin{leanitems}
\leanitem[1]{Part157Hell,Part157Repl,Part156Master}; \href{\leanrepo}{\leanstepC};
\href{\leanrepo}{\leanstepF}.
In Lean, $\bar\Delta$ is \texttt{HbarW}.
\leanitem[2]{Part157Gp,Part157Mgf,Part158Final,Part156Comp}. Lean proves
Lemma~\ref{lem:gbl} as \texttt{gaussGP}, by Ornstein--Uhlenbeck
interpolation.
\leanitem[3]{Part159Hid,Part159Uncond,Part227Frame,Part236Frame,Part290Frame,Part293World}
for $L_W$;
\leanlist{Part160Split,Part188Poincare,Part151Poinc,Part182Schur,Part291WeightedFrame,Part293FrameAdapter}
for $L_A$; \href{\leanrepo}{\leanstepC}, \href{\leanrepo}{\leanstepD},
\href{\leanrepo}{\leanstepJ}.
\end{leanitems}

\paragraph{Assembly.}
Say that $\Est$ \emph{errs} if it does not succeed. Following the chain~\eqref{eq:chain} of the
High-level overview, split the event that $\Est$ errs into three events:
$\Dd\notin\Gg$; $\Dd\in\Gg$ and $|\sqrt R-\sqrt{\tilde q(\Dd)}|>\Delta_s$; or neither
occurs and $\Est$ still errs. By Lemma~\ref{lem:dw}, the first
event has probability at most
$\delta_{\mathrm{hid}}+\mu_W/\Delta_W^2+\delta_{\mathrm{ext}}$ over $\Dd$.
By Lemma~\ref{lem:chernoff}, the second event has probability at most
$2e^{-K\Delta_s^2}$ over $r$. In the third event, let
$Q\coloneqq\EE_{\Omega_A}\tilde p$. Lemmas~\ref{lem:da}
and~\ref{lem:chernoff} give $|\sqrt R-\sqrt Q|\le\Delta_A+\Delta_s$, and Lemma~\ref{lem:dw}
gives $|\sqrt{\EE_{\Omega_W}p}-\sqrt{\tilde p}|\le\Delta_W$. Let
$\hat z_+\coloneqq\max(1/100,\Gamma\max(0,\sqrt Q-\Delta_A-\Delta_s)^2)$ and
$\hat z_-\coloneqq\Gamma(\sqrt Q+\Delta_A+\Delta_s)^2/10$. By the definition of $U$, as in the proof of
Lemma~\ref{lem:window}, $\Est$ then errs only if $p>\hat z_+$, or if $p<\hat z_-$ and
$\hat z_->1/1000$. These are tail events of the two links marked ``conc.'' in~\eqref{eq:chain}. Fix reals
$\sigma_+>0>\sigma_-$. Let $E_\pm$ be explicit functions of
$(\Delta_A+\Delta_s,\Delta_W,\vartheta_W,\vartheta_A,\sigma_\pm)$ and of $\Gamma$, with no further dependence on $n$,
and let $\Bb_\pm\coloneqq e^{nE_\pm}$. Only the Lean proof states $E_\pm$. Assume that
$E_+\le0$, $E_-\le0$, $\Delta_A+\Delta_s+\Delta_W<1/\sqrt{100\Gamma}$ and
$2(\Delta_A+\Delta_s)\le1/\sqrt{100\Gamma}$. For each $\Dd\in\Gg$, over $(\Omega_W,\Omega_A)$, Lemmas~\ref{lem:conc}
and~\ref{lem:frame}, with the bound of Lemma~\ref{lem:dw}, give
probability at most $\Bb_+$ that $p>\hat z_+$, and at most
$\Bb_-$ that $p<\hat z_-$ with $\hat z_->1/1000$. Hence the probability that $\Est$ errs is at most
\[\delta(n)\coloneqq \delta_{\mathrm{hid}}+\mu_W/\Delta_W^2+\delta_{\mathrm{ext}}
+\Bb_++\Bb_-+2e^{-K\Delta_s^2}~\leanmark{1}.\]
Lemma~\ref{lem:numcert} below will show that, for every integer
$n\ge8{,}000$, $\Est$ errs with probability at most $0.01$: by
$\delta(n)\le0.01$ if $n\ge11{,}883$, and by a finer form of $\delta(n)$ if
$n\le11{,}882$. Hence
$\PP\bigl(p\le U\le\max(10p,1/100)\bigr)\ge0.99$ for every
such $n$ and every $t\in\RR$, which is the success bound of the \bst{}.
\begin{leanitems}
\leanitem*[1]{\leanstepF}. In Lean, the bound is \texttt{network\_master}
in \texttt{Part196Master}.
\end{leanitems}

\paragraph{Numerical certificates.}
\begin{lemma}\label{lem:numcert}
  For every integer $n\ge8{,}000$ and every $t\in\RR$, $\Est$ errs with
  probability at most $0.01$. If $n\ge11{,}883$, Lean gives explicit
  constants, which depend on $n$, that satisfy the assumptions of the
  Assembly and give $\delta(n)\le0.01$. If $n\le11{,}882$, the
  probability is at most $\delta_{\mathrm{low}}<0.00874$, for
  $\delta_{\mathrm{low}}$ defined below.
\end{lemma}
Lean proves Lemma~\ref{lem:numcert} on four ranges of $n$:
$n\ge 16{,}200$~\leanmark{1}; $12{,}225\le n\le16{,}199$~\leanmark{2};
$11{,}883\le n\le 12{,}224$~\leanmark{3}; and $8{,}000\le n\le 11{,}882$~\leanmark{4}.
Lean checks the numerical certificates by exact rational enclosures,
including Taylor bounds for $\exp$. It splits the ranges into sub-ranges,
some of a single $n$, with their own constants. When a tail certificate
covers several values of $n$ with fixed exponent bounds, Lean checks it at
the least of them and transfers it to larger $n$, since
$E_\pm\le0$ makes $\Bb_\pm$ nonincreasing in $n$. On the first three ranges, $\Delta_s=10^{-5}$, so
$2e^{-K\Delta_s^2}=2e^{-100}$.

For example, for $16{,}384\le n\le16{,}500$, Lean uses the constants
\[\begin{aligned}
&\Delta_W=0.0024160420,&&\Delta_A=0.006583+10^{-20},&&\Delta_s=10^{-5},\\
&\vartheta_W=1.377,&&\vartheta_A=2.036,&&\sigma_+=0.1528,\quad\sigma_-=-0.4038,
\end{aligned}\]
and checks
\[\begin{aligned}
&E_+\le-0.000364553074,&&E_-\le-0.0003551354576,\\
&\Bb_++\Bb_-\le0.005519,&&\delta_{\mathrm{hid}}+\mu_W/\Delta_W^2+\delta_{\mathrm{ext}}\le0.00448,
\end{aligned}\]
so that $\delta(n)\le0.00448+0.005519+2e^{-100}<0.01$~\leanmark{5}.

On the last range, Lean uses a finer failure bound. Let $\Gg_{\mathrm{low}}$
be the intersection of the six events on $\Dd$ that the proof
for this range uses; it differs from $\Gg$. On $\Gg_{\mathrm{low}}$, Lean uses the single values
$\vartheta_A=1.56$, $\Delta_A=0.0064$ and $\Delta_s=5\cdot10^{-6}$. In place
of the one threshold $\Delta_W$ on $\bar\Delta$ and the one constant
$\vartheta_W$, it uses two nonnegative scalars $\varpi_1,\varpi_2$, which only the Lean
proof defines, with $\varpi\coloneqq \varpi_1+\varpi_2\le0.324425$ and $L_W\le(0.419+\varpi)/n$. Lean
covers the values of $\bar\Delta$ and $\varpi$ by finite intervals. On each pair
of intervals, it checks at $n=8{,}000$ a bound on $\min(1,\Bb_++\Bb_-)$, with
finer exponents in place of $E_\pm$; the same bound holds for every
$n\ge8{,}000$. Beyond the last interval for $\bar\Delta$, it uses the bound $1$.
It averages these bounds over $\Gg_{\mathrm{low}}$ through moments of
$\bar\Delta^2$, $\varpi_1^2$ and $\varpi_2$. The probability that $\Est$ errs is then at
most
\[\delta_{\mathrm{low}}\coloneqq0.0026872+0.0057+0.0003431447924852+10^{-9}
=0.0087303457924852.\]
The first term bounds $\PP(\Dd\notin\Gg_{\mathrm{low}})$. The second and
third terms together bound the averaged interval bounds, including the bound
$1$ beyond the last interval. The last term bounds $2e^{-K\Delta_s^2}=2e^{-25}$~\leanmark{6}. On $\Gg_{\mathrm{low}}$, the frame bounds of
Lemma~\ref{lem:frame} hold with $\vartheta_W=0.743425$ and $\vartheta_A=1.56$.

For odd $n\ge8{,}001$, the values $\bar h_{g,i}$ and $\bar h^{(2)}_{g,i}$ of the
finite rule of step~5 are within $10^{-90}n^{-10}a$ of $\EE_VX_{g,i}$ and within
$3\cdot10^{-90}n^{-10}a^2$ of $\EE_VX_{g,i}^2$, for every real $a$ with
$|s_g\alpha_i|\le a$ and $0\le\sqrt{1-s_g^2}\,\alpha^\perp_i\le a$. For every odd $n\ge8{,}001$, on the good event of the proof for the range of
$n$, Lean turns these errors into a change of at most $10^{-20}$ in
$\sqrt{\tilde q}$, the fourth term in the proof of Lemma~\ref{lem:da}; for the
example above, it is the term $10^{-20}$ in $\Delta_A$~\leanmark{7}.
\begin{leanitems}
\leanitem*[1]{\leanstepG}.
\leanitem*[2]{\leanstepH}.
\leanitem*[3]{\leanstepI}.
\leanitem*[4]{\leanstepJ}.
\leanitem[5]{Part186GLedC,Part186GCombC,Part160Ledger2Gen,Part156Num}.
\leanitem[6]{Part293Ledger,Part293Assembly,Part293EndpointInterface,Part236Master}.
\leanitem[7]{Part278Err,Part278StBase3,Part281Odd,Part293Arithmetic}.
\end{leanitems}

\paragraph{Total cost.}
The step costs sum to $2m_W^2+1024n^5+4m_WJ_W+O(n^{4.5})+2Km_F^2+O(Knm_F)$,
where $1024n^5$ is the cost of step~5 and $4m_WJ_W$ that of the rule
$B_{J_W}$ of step~2; for $8{,}000\le n\le16{,}384$ the first is $22$ to $27$
times the second. The leading term
$2m_W^2$ is the cost of the product $W_2\psi$ of step~2; one exact sample
computes the product $W_2\relu(W_1u)$, which has the same leading cost. The $K$ evaluations of $F$ cost about
$2Km_F^2\approx10^{24}$ at $n=8{,}000$, far below $2n^7>4\cdot10^{27}$. For
every $n\ge8{,}000$, the total is below $1.0003\cdot2n^7$~\leanmark{1}.
\begin{leanitems}
\leanitem[1]{Part282Total,Part282JoinCost,Part279ScheduleOut,Part289Cost}; \href{\leanrepo}{\leanstepK}.
\end{leanitems}

The cost model in Lean charges nothing for rational constants or for reading
the seed. The Cost items of steps~1, 2, 3, 5 and~6 estimated this uncharged
work. The sum of these estimates is below $2\cdot10^{18}$ at $n=8{,}000$,
less than $10^{-9}$ times $2n^7$. Almost all of the sum comes from the
Gaussian draws of step~6. These estimates are not checked in Lean.

\section{LLM usage}\label{sec:llm-usage}
LLMs from multiple providers was used to discover, prove, and explain this result, via the following methodology.

\paragraph{Discovery via informal proof.}
LLMs produced the proofs under human direction. Claude Opus~5.5
(medium effort) coordinated the work, wrote and challenged paper proofs, and
checked each result before the result was recorded. GPT-6.1-Sol (Codex, medium
effort) challenged paper proofs. GPT-6-Astra (Codex, medium effort) wrote the
proof for $8{,}000\le n\le11{,}882$. By volume of written reports the paper
work was split roughly $57\%$ Codex and $42\%$ Opus. Each claim was challenged
by a second model before the claim was relied on.

\paragraph{Formal proof.}
LLMs formalized the proofs in Lean. GPT-6.1-Sol formalized most of the
Lean development, including the shared library and the certificates for
$n\ge12{,}225$. GPT-6-Astra (high effort) formalized the proof for
$8{,}000\le n\le11{,}882$, the band-wise root certificates and the final
assembly. Claude Opus~5.5 formalized the treatment of odd dimensions, the
operation count, the sharp root bounds and several assemblies. Claude
Sonnet~5.5 (medium effort) formalized some numerical tables. Of the $202$ Lean
parts in the import closure of the two final theorems, roughly $77\%$ are Sol's, $16\%$ Opus's, $6\%$
Astra's and $1\%$ Sonnet's. The Lean kernel checks the proof of the final
statement. A Lean theorem can compile even when its statement is not the
intended claim \citep[cf.][Sections~1 and~4]{ilin-nugent-2026}; for that reason GPT-6-Astra (high
effort) and Claude Opus~5.5 separately compared this note's text with the Lean
statements of the setting, of $\Est$ and of the conclusion.

\paragraph{Explanation.}
The human author iteratively refined a rule-set for Claude Opus~5.5 improving 
the exposition, and made occasional direct changes, until the logic flow of 
the argument became sufficiently organized and clear to the author.
 

\section{Related work}

\paragraph{ARC's program.}
The Alignment Research Center (ARC) proposes to estimate the probability that
a catastrophe detector fires on a model's outputs, by analyzing the model,
because sampling is slow when failures are rare
\citep{arc-2025-competing,arc-2026-mikes-eye}: $\Kk$ samples are unlikely to
show a failure much rarer than $1/\Kk$, and so cannot estimate its rate to
within a constant factor. The network $A$, the detector $C$ and the distribution of $x$ play the
roles of ARC's model, detector and input distribution, and the failure rate
$p$ of the \bst{} is the quantity ARC's scheme needs. ARC's
\emph{Matching Sampling Principle} is Conjecture~\ref{conj:msp} below. Let
$M_\theta(c,x)$ be a model with weights $\theta$, a random context $c$ and a
random input $x$; in this paragraph and the next, $c$ and $x$ are ARC's.
\begin{conjecture}[Matching Sampling Principle, ARC]\label{conj:msp}
For every architecture $M$ there is an estimator with the following
property. For every $\theta$ there is an explanation $\pi$ of size
$O(|\theta|)$, obtained in some unspecified way, such that for every
tolerance $\varepsilon>0$ the estimator uses $\theta$, $\pi$, $c$ and
$\varepsilon$ to estimate the mean output $\EE_x[M_\theta(c,x)]$, and:
averaged over the context $c$, its mean squared error is at most
$\varepsilon^2$ times the average variance $\EE_c\operatorname{Var}_x
M_\theta(c,x)$; it runs in time $O(\varepsilon^{-2})$ forward passes of
$M_\theta$; and it is \emph{mechanistic}, a requirement that ARC leaves
informal, i.e., it estimates deductively from the structure of $M_\theta$
instead of from sampled behavior.
\end{conjecture}
Averaging $\varepsilon^{-2}$ independent samples of
$M_\theta(c,x)$ achieves exactly this accuracy at exactly this cost, but is
not mechanistic; this is why the principle speaks of \emph{matching} sampling. Random weights can be
placed in the context; ARC's results for random networks take this form and
need no explanation, as the paragraph ``Mechanistic estimation for random
networks'' will discuss.

Now let $M_\theta(c,x)$ be a failure indicator, equal to $1$ when the
detector fires and $0$ otherwise, so that its mean is the failure rate $p$.
A $0/1$ variable with mean $p$ has variance $p(1-p)$, which is about $p$ when
failures are rare. As a benchmark, consider a single context with failure
rate $p$ and an error of about $\varepsilon\sqrt p$. This error is smaller
than $p$ itself only when $\varepsilon^2\lesssim p$, so that the estimate is
right to within a constant factor, and then the time allowed is about $1/p$
forward passes. That is also what sampling costs to estimate a probability
$p$ to within a constant factor. So Conjecture~\ref{conj:msp}, as stated, does not
promise that a mechanistic estimator beats sampling on rare failures; ARC
hopes to do better on structured models. The program $\Est$ of the \bst{} beats
naive-99\% samplers on the target of $\Est$, $p\le U\le\max(10p,1/100)$: for every
$p\le1/1000$, the cost of $\Est$ is more than $447$ times smaller than the expected cost
of any naive-99\% sampler (Section~\ref{sec:comparison}). We prove no such lower bound for samplers that
need to succeed only for most networks, nor for importance sampling or
splitting.

\paragraph{Mechanistic estimation for random networks.}
\citet{arc-2026-wide-mlps}
estimate the expected output of a wide random MLP under Gaussian input by
propagating cumulants, without running the network. Their theorems
assume fixed depth and polynomial activations as the width $n\to\infty$. For
tolerances $\varepsilon$ at least inverse-polynomial in $n$, they give mean
squared error $O(\varepsilon^2)$ averaged over the weights. When
$\varepsilon^2=O(1/n^2)$, their running time is $O(n/\varepsilon^2)$, against
$\Theta(n^2/\varepsilon^2)$ for Monte Carlo. Their experiments use ReLU at
width $256$ and beat sampling in error per FLOP, including on rare thresholded events. Whether such methods can match sampling
at large depth is open there, and ARC's White-Box Estimation
Challenge \citep{arc-2026-challenge,aicrowd-2026-phase2} asks contestants to estimate the same mean output, at width $256$ and depth
$32$ in Phase~1 and width $1{,}024$ and depth $16$ in Phase~2, without
requiring mechanistic methods. For one layer, \citet{hilton-2025-halfspaces} deductively
estimates the probability that all output neurons are active (an orthant
probability), matching the mean squared error of sampling averaged over the
weights, up to logarithmic factors in the running time. Compared with this
work, the \bst{} covers only depth two and prescribed width scalings,
but it concerns a probability over a non-Gaussian input distribution produced
by a larger network, holds with explicit constants at every $n\ge8{,}000$,
and is checked in Lean.

\paragraph{Other approaches.}
Rare-event methods estimate failure probabilities of networks with fewer
samples, by importance sampling or multilevel
splitting \citep{webb-2019,wu-hilton-2025-rare}, but still evaluate the network many
times and give no guarantee of the kind proved here. Moment propagation
through ReLU layers under a Gaussian approximation is standard in Bayesian
deep learning \citep{wu-2019-dvi}, with empirical checks but no error bounds. Sound probabilistic
verifiers bound output probabilities rigorously by branch and bound over the
input space \citep{boetius-2025}, for a given network, trained ones included;
the experiments of \citet{boetius-2025} use input dimensions far smaller than
$2n\ge16{,}000$, and that paper gives no cost guarantee comparable to
the \bst{}. Infinite-width theory \citep{matthews-2018,lee-2018-nngp,hanin-2023}
describes random networks over the random weights, whereas $\Est$
estimates a probability over inputs for one draw of the weights. \citet{gross-2024-compact} prove compact accuracy bounds for small trained one-layer attention-only
transformers, with proofs shorter than brute-force enumeration; our result
instead bounds the cost of an estimator, for random networks at frontier
input widths.

\begin{thebibliography}{99}\small
\bibitem[AIcrowd(2026)]{aicrowd-2026-phase2} AIcrowd. \emph{Phase 2 of the ARC White-Box Estimation Challenge is live.} AIcrowd forum, August 23, 2026.
\bibitem[Bayati and Montanari(2011)]{bayati-montanari-2011} M. Bayati, A. Montanari. \emph{The dynamics of message passing on dense graphs, with applications to compressed sensing.} IEEE Trans.\ Inform.\ Theory 57(2):764--785, 2011. arXiv:1001.3448.
\bibitem[Bennett et~al.(2008)]{bcct-2008} J. Bennett, A. Carbery, M. Christ, T. Tao. \emph{The Brascamp--Lieb inequalities: finiteness, structure and extremals.} Geom.\ Funct.\ Anal.\ 17(5):1343--1415, 2008.
\bibitem[Boetius et~al.(2025)]{boetius-2025} D. Boetius, S. Leue, T. Sutter. \emph{Solving probabilistic verification problems of neural networks using branch and bound.} ICML 2025.
\bibitem[Bolthausen(2014)]{bolthausen-2014} E. Bolthausen. \emph{An iterative construction of solutions of the TAP equations for the Sherrington--Kirkpatrick model.} Comm.\ Math.\ Phys.\ 325:333--366, 2014.
\bibitem[Boucheron et~al.(2013)]{blm-2013} S. Boucheron, G. Lugosi, P. Massart. \emph{Concentration Inequalities: A Nonasymptotic Theory of Independence.} Oxford University Press, 2013.
\bibitem[Cho and Saul(2009)]{cho-saul-2009} Y. Cho, L. K. Saul. \emph{Kernel methods for deep learning.} Advances in Neural Information Processing Systems 22 (NeurIPS 2009).
\bibitem[DeepSeek-AI(2026)]{deepseek-2026-v4} DeepSeek-AI. \emph{DeepSeek-V4-Pro}, model card and configuration. Hugging Face, 2026. \url{https://huggingface.co/deepseek-ai/DeepSeek-V4-Pro}
\bibitem[DLMF(2026)]{dlmf} F. W. J. Olver et~al.\ (eds.). \emph{NIST Digital Library of Mathematical Functions}, Section~5.12, Release~1.2.8 of 2026-09-15. \url{https://dlmf.nist.gov/5.12}
\bibitem[Devroye et~al.(2022)]{dmr-2022} L. Devroye, A. Mehrabian, T. Reddad. \emph{The total variation distance between high-dimensional Gaussians with the same mean.} arXiv:1810.08693v6, 2022.
\bibitem[Gross et~al.(2024)]{gross-2024-compact} J. Gross, R. Agrawal, T. Kwa, E. Ong, C. H. Yip, A. Gibson, S. Noubir, L. Chan. \emph{Compact proofs of model performance via mechanistic interpretability.} NeurIPS 2024.
\bibitem[Halmos and Sunder(1978)]{halmos-sunder-1978} P. R. Halmos, V. S. Sunder. \emph{Bounded Integral Operators on $L^2$ Spaces.} Ergebnisse der Mathematik und ihrer Grenzgebiete 96, Springer, 1978.
\bibitem[Hanin(2023)]{hanin-2023} B. Hanin. \emph{Random neural networks in the infinite width limit as Gaussian processes.} Ann.\ Appl.\ Probab.\ 33(6A):4798--4819, 2023.
\bibitem[Hilton(2025)]{hilton-2025-halfspaces} J. Hilton. \emph{Deductive estimation for random halfspace intersections.} Note dated May 2, 2025, published by the Alignment Research Center in May 2026.
\bibitem[Hilton et~al.(2026)]{arc-2026-challenge} J. Hilton, P. Christiano, W. Wu. \emph{Announcing the ARC White-Box Estimation Challenge.} Alignment Research Center blog, June 2, 2026.
\bibitem[Ilin and Nugent(2026)]{ilin-nugent-2026} V. Ilin, B. Nugent. \emph{Sorries are not the hard part: an expert-review case study of a semi-autonomous formalization.} arXiv:2606.13925, 2026.
\bibitem[Lee et~al.(2018)]{lee-2018-nngp} J. Lee, Y. Bahri, R. Novak, S. S. Schoenholz, J. Pennington, J. Sohl-Dickstein. \emph{Deep neural networks as Gaussian processes.} ICLR 2018.
\bibitem[Lehec(2014)]{lehec-2014} J. Lehec. \emph{Short probabilistic proof of the Brascamp--Lieb and Barthe theorems.} Canad.\ Math.\ Bull.\ 57(3):585--597, 2014.
\bibitem[Matthews et~al.(2018)]{matthews-2018} A. G. de G. Matthews, J. Hron, M. Rowland, R. E. Turner, Z. Ghahramani. \emph{Gaussian process behaviour in wide deep neural networks.} ICLR 2018.
\bibitem[Meta(2026)]{meta-2026-muse} Meta. \emph{Introducing Muse Spark Open.} Meta AI blog, October 15, 2026, \url{https://research.meta.ai/blog/muse-spark-open}; architecture from Hugging Face, \emph{Welcome Muse Spark Open}, Hugging Face blog, 2026, \url{https://huggingface.co/blog/meta-muse-spark-open}; configuration from the public copy \url{https://huggingface.co/unsloth/Muse-Spark-Open-816B-A42B} and the class \texttt{MuseSparkMoE} in Hugging Face Transformers.
\bibitem[Moonshot AI(2026)]{kimi-2026-k3} Moonshot AI. \emph{Kimi K3}, model card and configuration. Hugging Face, 2026. \url{https://huggingface.co/moonshotai/Kimi-K3}
\bibitem[Neyman et~al.(2025)]{arc-2025-competing} E. Neyman, V. Lecomte, W. Wu, M. Winer, J. Hilton, G. Robinson. \emph{Competing with sampling.} Alignment Research Center blog, November 18, 2025.
\bibitem[Okamoto(1958)]{okamoto-1958} M. Okamoto. \emph{Some inequalities relating to the partial sum of binomial probabilities.} Ann.\ Inst.\ Statist.\ Math.\ 10(1):29--35, 1958.
\bibitem[Qwen Team(2026)]{qwen-2026-38} Qwen Team, Alibaba. \emph{Qwen3.8-2.4T-A95B}, model card and configuration. Hugging Face, 2026. \url{https://huggingface.co/Qwen/Qwen3.8-2.4T-A95B}
\bibitem[Vershynin(2018)]{vershynin-2018-hdp} R. Vershynin. \emph{High-Dimensional Probability: An Introduction with Applications in Data Science.} Cambridge University Press, 2018.
\bibitem[Wald(1945)]{wald-1945} A. Wald. \emph{Sequential tests of statistical hypotheses.} Ann.\ Math.\ Statist.\ 16(2):117--186, 1945.
\bibitem[Webb et~al.(2019)]{webb-2019} S. Webb, T. Rainforth, Y. W. Teh, M. P. Kumar. \emph{A statistical approach to assessing neural network robustness.} ICLR 2019.
\bibitem[Winer(2026)]{arc-2026-mikes-eye} M. Winer. \emph{A Mike's-eye view of ARC's research.} Alignment Research Center blog, June 9, 2026.
\bibitem[Wu and Hilton(2025)]{wu-hilton-2025-rare} G. Wu, J. Hilton. \emph{Estimating the probabilities of rare outputs in language models.} ICLR 2025.
\bibitem[Wu et~al.(2019)]{wu-2019-dvi} A. Wu, S. Nowozin, E. Meeds, R. E. Turner, J. M. Hern\'andez-Lobato, A. L. Gaunt. \emph{Deterministic variational inference for robust Bayesian neural networks.} ICLR 2019.
\bibitem[Wu et~al.(2026)]{arc-2026-wide-mlps} W. Wu, V. Lecomte, M. Winer, G. Robinson, J. Hilton, P. Christiano. \emph{Estimating the expected output of wide random MLPs more efficiently than sampling.} arXiv:2605.05179, 2026.
\bibitem[Yang(2019)]{yang-2019-tp1} G. Yang. \emph{Tensor Programs I: Wide feedforward or recurrent neural networks of any architecture are Gaussian processes.} NeurIPS 2019. Full version: arXiv:1910.12478v3, 2021.
\bibitem[Z.ai(2026)]{zai-2026-glm53} Z.ai. \emph{GLM-5.3}, model card and configuration. Hugging Face, 2026. \url{https://huggingface.co/zai-org/GLM-5.3}
\end{thebibliography}

\appendix
\section{Network sizes}\label{app:sizes}

Table~\ref{tab:sizes-app} repeats Table~\ref{tab:sizes}, with fuller notes in
its caption. The input dimension of $A$ is comparable to the
embedding widths of the models in Table~\ref{tab:sizes-app}, but its hidden width and weight count are
far larger: $A$ has about $n^4$ weights, $4.1\cdot10^{15}$ at $n=8{,}000$,
while all of Kimi K3 has a reported $2.8\cdot10^{12}$
parameters~\citep{kimi-2026-k3}. The failure metric $F$ has hidden width
$\lceil n^{1.5}\rceil=715{,}542$ at $n=8{,}000$, which is closer in size to one
mixture-of-experts layer, i.e., to the values $524{,}288\le h\le2{,}752{,}512$
of Table~\ref{tab:sizes-app}. The transformer MLPs also differ in kind: they are
trained, use gated activations rather than ReLU, have one
hidden layer, and sit between attention layers, while $A$ is a random dense
ReLU network with two hidden layers.

\begin{table}[h]
\centering\small
\sizestable
\caption{The agent and world networks of the \bst{} at $n=8{,}000$, and the MLP blocks of large open-weight transformers, read
from their published configurations (October 2026). The last column counts
the parameters of the whole model: the weights of $A$ and of $W$, and the totals stated
in each model card or release announcement; GLM-5.3's model card states no
total, so its entry is the parameter count that Hugging Face reports for its checkpoint. Kimi K3's routed experts act on a $3{,}584$-dimensional
projection of its $7{,}168$-dimensional embedding, and Muse Spark Open's on a $4{,}096$-dimensional projection of its $8{,}192$-dimensional embedding. The mixture-of-experts
layers have $896$ routed experts of width $3{,}072$ with $16$ used per token
(Kimi K3), $512$ of width $2{,}048$ with $10$ per token (Qwen3.8), $384$ of
width $3{,}072$ with $6$ per token (DeepSeek V4 Pro), $256$ of width $4{,}096$
with $8$ per token (Muse Spark Open) and $256$ of width $2{,}048$ with $8$ per
token (GLM-5.3). All but Muse Spark Open also have one or two shared experts.
The configuration of Muse Spark Open sets no shared experts. The official Muse Spark Open checkpoint is gated, so its figures are from the configuration of a public copy of the checkpoint and from the Transformers implementation.}
\label{tab:sizes-app}
\end{table}

\section{Lean files}\label{app:files}
The extracted project has $5{,}833$ modules: the import closure of the modules of the two Lean theorems that state the bounds of the \bst{}, \texttt{Part293.classifier\_stated2\_8000} and \texttt{Part289.progOut\_run2}, together with two modules, \texttt{Part288Audit} and \texttt{Part289All}, that contain only imports and axiom-audit commands. Of these modules, $4{,}290$ define a declaration that the proofs of these two theorems use. The table assigns every module to one group (groups S0--S10 follow Section~\ref{sec:program}: S0 is the setting, S1 the definition of $\Est$, S2 and S3 what is known after step~4, S4 what is known after step~6, S5 the assembly, S6--S9 the numerical certificates on the four ranges of $n$, and S10 the total cost; S11 holds modules that are compiled because they are imported, but contribute no declaration to these proofs). A script checks that the groups cover every module. Parts are module-name prefixes \texttt{BstLean.Part<k>}.
{\footnotesize
\begin{longtable}{@{}l l p{0.30\textwidth} r r p{0.34\textwidth}@{}}
\toprule Step & Group & Parts & Files & Used & Role \\ \midrule \endhead
S11 & G00 & \raggedright 29, 33, 35--46, 48--53, 56--64, 66, 69, 71--74, 78, 80--82, 84, 86, 88--90, 93, 105--109, 112, 130--135, 138--145, 149--154, 156--180, 182, 185--188, 192--194, 197, 199--201, 204--205, 208--213, 216--217, 220--223, 235--236, 238--239, 241--245, 247, 249, 251--257, 260--262, 264--266, 269, 271, 273, 275--276, 278--279, 281--282, 287--290, 292--293 & 1543 & 0 & \raggedright These imported modules contribute no declaration to either final theorem. \tabularnewline
S0 & G01 & \raggedright 130, 132, 156--157 & 11 & 11 & \raggedright The definitions specify the Gaussian networks, normalized input, target probability, and success event. \tabularnewline
S0 & G02 & \raggedright 29, 33, 35--36, 41, 48, 53, 56, 58, 73, 80, 105--108, 131, 135, 139, 144, 146 & 41 & 41 & \raggedright Gaussian and spherical integration supplies moments, symmetry, measurability, and polar coordinates. \tabularnewline
S0 & G03 & \raggedright 130, 132--133, 142, 145 & 24 & 24 & \raggedright The ReLU kernel and layer moment bounds control the means and covariances of hidden vectors. \tabularnewline
S0 & G04 & \raggedright 130, 140, CostModel & 3 & 3 & \raggedright The program language charges one operation for each arithmetic operation, square root, and ReLU evaluation, and counts the operations of matrix-vector products. \tabularnewline
S1 & G05 & \raggedright 157, 169, 272 & 14 & 14 & \raggedright The program computes the two centres and the height tables and draws the Gaussian mixtures; the closed forms of the height moments are proved for even dimensions. \tabularnewline
S1 & G06 & \raggedright 141, 150, 159, 201 & 31 & 31 & \raggedright Finite binomial sums and truncated series approximate the Gaussian ReLU means and the spherical constants. \tabularnewline
S1 & G07 & \raggedright 278 & 11 & 11 & \raggedright The arithmetic rule for odd dimensions approximates the exact height moments and bounds the resulting table perturbation. \tabularnewline
S1 & G08 & \raggedright 289 & 2 & 2 & \raggedright The scale floor equals 0.17 below dimension 11{,}883 and 0.20 at or above it before multiplication by $2/n$. \tabularnewline
S2 & G09 & \raggedright 134, 157 & 5 & 5 & \raggedright Gaussian regression separates each last weight matrix into its retained action on a centre and an independent residual. \tabularnewline
S2 & G10 & \raggedright 137, 140, 143, 204, 210 & 23 & 23 & \raggedright The computed World centre and residual scale approximate the corresponding input averages. \tabularnewline
S2 & G11 & \raggedright 138, 157, 159 & 18 & 18 & \raggedright Gaussian Hellinger bounds transfer the World centre and scale errors to square roots of event probabilities. \tabularnewline
S2 & G12 & \raggedright 205, 214, 218 & 68 & 68 & \raggedright Conditional kernel expansions bound the World mean, covariance, direction error, and hidden energy fluctuations. \tabularnewline
S2 & G13 & \raggedright 227, 233 & 22 & 22 & \raggedright World event bounds control the centre, noise window, frame, and integrated Hellinger error. \tabularnewline
S2 & G14 & \raggedright 232, 236 & 14 & 14 & \raggedright Integration over radius and frame intervals converts radius and frame moments into a bound on the classifier failure probability. \tabularnewline
S3 & G15 & \raggedright 107, 109, 148, 151--152, 244, 246 & 35 & 35 & \raggedright Gaussian concentration, Poincaré inequalities, and matrix tail bounds control norms and quadratic forms. \tabularnewline
S3 & G16 & \raggedright 105, 112, 134, 136 & 14 & 14 & \raggedright Gaussian correlation estimates and second-moment identities bound the operator norm of averaged residual directions. \tabularnewline
S3 & G17 & \raggedright 157--158 & 14 & 14 & \raggedright The Gaussian Brascamp--Lieb inequality converts a frame bound into an exponential-moment bound for an event probability. \tabularnewline
S3 & G18 & \raggedright 151, 153, 229 & 42 & 42 & \raggedright Jacobian and inverse-radius estimates give covariance and moment bounds for the Agent hidden vector. \tabularnewline
S3 & G19 & \raggedright 149, 154, 160, 194, 198, 206, 211 & 86 & 86 & \raggedright Mean and deviation estimates bound the error and energy of the computed Agent centre. \tabularnewline
S3 & G20 & \raggedright 160, 195, 215, 222 & 47 & 47 & \raggedright Covariance, centre error, and residual-energy tails combine into the Agent frame bound. \tabularnewline
S3 & G21 & \raggedright 188 & 9 & 9 & \raggedright A capped reciprocal weight and the Poincaré inequality control the average of squared residual directions. \tabularnewline
S3 & G22 & \raggedright 199, 208, 235, 238 & 40 & 40 & \raggedright Schur bounds split input pairs into central, intermediate, and exceptional regions to control the Agent covariance. \tabularnewline
S3 & G23 & \raggedright 200, 209, 217, 221, 245, 273 & 50 & 50 & \raggedright Conditional Gaussian pair estimates control the tangent Jacobian and its error after the first Agent layer. \tabularnewline
S3 & G24 & \raggedright 219 & 22 & 22 & \raggedright Rotation and moment bounds control the input gradient of the hidden vector's coordinate along the computed Agent centre. \tabularnewline
S4 & G25 & \raggedright 161--165 & 45 & 45 & \raggedright Hellinger and chi-square bounds compare location-scale mixtures and separate fluctuation from bias. \tabularnewline
S4 & G26 & \raggedright 166--168, 212 & 56 & 56 & \raggedright Height-conditioned Jacobian moments bound the variation of the Agent location and residual scale. \tabularnewline
S4 & G27 & \raggedright 169--171, 202, 213, 216, 220, 223, 239 & 130 & 130 & \raggedright Moment bounds for the finite height tables control interpolation error, location bias, and scale bias. \tabularnewline
S4 & G28 & \raggedright 228, 241 & 38 & 38 & \raggedright Gaussian height tails and variable bias parameters bound the Hellinger error of Agent replacement. \tabularnewline
S4 & G29 & \raggedright 207, 242, 248 & 35 & 35 & \raggedright Conditional variance and reciprocal expansions reduce the Hellinger error of the Agent mixture. \tabularnewline
S4 & G30 & \raggedright 289 & 8 & 8 & \raggedright The bias and mixture estimates remain valid for the smaller positive scale floor. \tabularnewline
S5 & G31 & \raggedright 156--157, 196, 203 & 20 & 20 & \raggedright The two residual exponential-moment bounds, the replacement root bound, and the sampling error imply the required two-sided output guarantee. \tabularnewline
S5 & G32 & \raggedright 193, 197, 261 & 13 & 13 & \raggedright Exact interval arithmetic and finite series certify the scalar inequalities in the final probability bounds. \tabularnewline
S5 & G33 & \raggedright 192, 243 & 5 & 5 & \raggedright Parameterized World, Agent, and endpoint bounds connect the analytic estimates to the network classifier. \tabularnewline
S5 & G34 & \raggedright 293 & 2 & 2 & \raggedright The four dimension ranges combine into the success theorem for every integer $n\ge8{,}000$. \tabularnewline
S6 & G35 & \raggedright 172--181 & 142 & 142 & \raggedright Explicit bounds on exceptional probabilities and Hellinger errors establish the classifier guarantee at large dimensions without an upper dimension limit. \tabularnewline
S6 & G36 & \raggedright 182--188, 190 & 386 & 386 & \raggedright Sharper covariance, energy, and table bounds extend the upper-range guarantee down to dimension 16{,}200. \tabularnewline
S6 & G37 & \raggedright 278 & 45 & 45 & \raggedright The odd-dimension table perturbation transfers the upper-range guarantee to the stated arithmetic program. \tabularnewline
S7 & G38 & \raggedright 250, 260, 268 & 312 & 312 & \raggedright World event and Hellinger certificates supply the individual dimensions in the lower portion of the middle range. \tabularnewline
S7 & G39 & \raggedright 252, 257, 262, 266 & 352 & 352 & \raggedright Bounds on exceptional Agent events and height table errors control the Hellinger error of Agent replacement in the middle range. \tabularnewline
S7 & G40a & \raggedright 251, 253--254, 256, 258--259, 263 & 160 & 160 & \raggedright Local noise intervals supply centre, direction, energy, and covariance inputs for the middle-range frame bounds. \tabularnewline
S7 & G40b & \raggedright 255, 264 & 327 & 327 & \raggedright Checked Gaussian matrix and geometric bounds supply the local covariance and frame estimates below dimension 12{,}350. \tabularnewline
S7 & G40c & \raggedright 269 & 432 & 432 & \raggedright Local covariance, centre, and energy estimates bound the frame for dimensions 12{,}350 through 13{,}224. \tabularnewline
S7 & G41a & \raggedright 268 & 175 & 175 & \raggedright Finite noise-cell tables certify the scalar frame inequalities on dimensions 12{,}350 through 13{,}224. \tabularnewline
S7 & G41 & \raggedright 243, 249, 268 & 76 & 76 & \raggedright Rational endpoint certificates verify the scalar success inequalities for dimensions 12{,}225 through 13{,}224. \tabularnewline
S7 & G42 & \raggedright 267, 270, 278, 286 & 51 & 51 & \raggedright The row certificates give the stated-program guarantee on dimensions 12{,}225 through 13{,}199; group G45 assembles the rows through 13{,}224. \tabularnewline
S7 & G43 & \raggedright 271 & 35 & 35 & \raggedright Uniform World, centre, and moment bounds supply the upper portion of the middle dimension range. \tabularnewline
S7 & G44 & \raggedright 285 & 48 & 48 & \raggedright Interval estimates control the probability of a table error and the Hellinger error of Agent replacement in the upper middle range. \tabularnewline
S7 & G45 & \raggedright 287 & 24 & 24 & \raggedright The lower row certificates and upper interval certificates give contiguous coverage from dimension 12{,}225 to 16{,}199. \tabularnewline
S8 & G46 & \raggedright 275, 277 & 13 & 13 & \raggedright Parameterized frame and Hellinger bounds support the local estimates on dimensions 11{,}883 through 12{,}224. \tabularnewline
S8 & G47 & \raggedright 276, 280 & 83 & 83 & \raggedright Sharper table moments and exceptional-probability bounds control the Hellinger error of Agent replacement on dimensions 11{,}883 through 12{,}224. \tabularnewline
S8 & G48 & \raggedright 281 & 99 & 99 & \raggedright World certificates bound the hidden events and Hellinger error on dimensions 11{,}883 through 12{,}224. \tabularnewline
S8 & G49a & \raggedright 281 & 101 & 101 & \raggedright Rational cell and endpoint certificates verify the frame and failure bounds on dimensions 11{,}883 through 12{,}224. \tabularnewline
S8 & G49b & \raggedright 281 & 37 & 37 & \raggedright The World, Agent, and scalar bounds assemble the stated classifier guarantee on dimensions 11{,}883 through 12{,}224. \tabularnewline
S8 & G49 & \raggedright 281 & 85 & 85 & \raggedright Local centre, energy, and covariance bounds supply the frame estimates on dimensions 11{,}883 through 12{,}224. \tabularnewline
S9 & G50 & \raggedright 290, 293 & 17 & 17 & \raggedright Uniform World bounds control the noise window, exceptional probabilities, and Hellinger moments throughout the low band. \tabularnewline
S9 & G51 & \raggedright 291, 293 & 36 & 36 & \raggedright Uniform Agent centre and height-conditioned moments supply the location and residual-scale estimates on the low band. \tabularnewline
S9 & G52 & \raggedright 291, 293 & 99 & 99 & \raggedright Weighted Schur and Poincaré estimates with checked interval covers bound the Agent frame throughout the low band. \tabularnewline
S9 & G53 & \raggedright 292--293 & 67 & 67 & \raggedright A common event with bounded complement probability controls the stored location and energy errors throughout the low band. \tabularnewline
S9 & G54 & \raggedright 293 & 30 & 30 & \raggedright Bias and conditional-mixture estimates at the positive scale floor bound the Hellinger error of Agent replacement throughout the low band. \tabularnewline
S9 & G55 & \raggedright 293 & 13 & 13 & \raggedright A radius and frame cover at dimension 8{,}000 bounds the failure probability uniformly over the low band. \tabularnewline
S9 & G56 & \raggedright 293 & 11 & 11 & \raggedright The World, Agent, table, and frame estimates combine with the endpoint bound and exceptional-probability bounds to prove success throughout dimensions 8{,}000 through 11{,}882. \tabularnewline
S10 & G57 & \raggedright 189, 279 & 12 & 12 & \raggedright Explicit operation schedules bound the World pass, Agent tables, Gaussian-mixture draws, and final output. \tabularnewline
S10 & G58 & \raggedright 150, 282, 288 & 21 & 21 & \raggedright Counted arithmetic programs return the stated centres, tables, samples, and output and include the comparisons in the table search. \tabularnewline
S10 & G59 & \raggedright 289 & 3 & 3 & \raggedright The revised program returns the stated output and costs less than $1.0003\cdot2n^7$ operations for every $n\ge8{,}000$. \tabularnewline
\midrule
& Total & & 5833 & 4290 & \\
\bottomrule
\end{longtable}

}

\end{document}
